% Options for packages loaded elsewhere
\PassOptionsToPackage{unicode}{hyperref}
\PassOptionsToPackage{hyphens}{url}
\documentclass[
]{article}
\usepackage{xcolor}
\usepackage{amsmath,amssymb}
\setcounter{secnumdepth}{-\maxdimen} % remove section numbering
\usepackage{iftex}
\ifPDFTeX
  \usepackage[T1]{fontenc}
  \usepackage[utf8]{inputenc}
  \usepackage{textcomp} % provide euro and other symbols
\else % if luatex or xetex
  \usepackage{unicode-math} % this also loads fontspec
  \defaultfontfeatures{Scale=MatchLowercase}
  \defaultfontfeatures[\rmfamily]{Ligatures=TeX,Scale=1}
\fi
\usepackage{lmodern}
\ifPDFTeX\else
  % xetex/luatex font selection
\fi
% Use upquote if available, for straight quotes in verbatim environments
\IfFileExists{upquote.sty}{\usepackage{upquote}}{}
\IfFileExists{microtype.sty}{% use microtype if available
  \usepackage[]{microtype}
  \UseMicrotypeSet[protrusion]{basicmath} % disable protrusion for tt fonts
}{}
\makeatletter
\@ifundefined{KOMAClassName}{% if non-KOMA class
  \IfFileExists{parskip.sty}{%
    \usepackage{parskip}
  }{% else
    \setlength{\parindent}{0pt}
    \setlength{\parskip}{6pt plus 2pt minus 1pt}}
}{% if KOMA class
  \KOMAoptions{parskip=half}}
\makeatother
\usepackage{longtable,booktabs,array}
\newcounter{none} % for unnumbered tables
\usepackage{calc} % for calculating minipage widths
% Correct order of tables after \paragraph or \subparagraph
\usepackage{etoolbox}
\makeatletter
\patchcmd\longtable{\par}{\if@noskipsec\mbox{}\fi\par}{}{}
\makeatother
% Allow footnotes in longtable head/foot
\IfFileExists{footnotehyper.sty}{\usepackage{footnotehyper}}{\usepackage{footnote}}
\makesavenoteenv{longtable}
\usepackage{graphicx}
\makeatletter
\newsavebox\pandoc@box
\newcommand*\pandocbounded[1]{% scales image to fit in text height/width
  \sbox\pandoc@box{#1}%
  \Gscale@div\@tempa{\textheight}{\dimexpr\ht\pandoc@box+\dp\pandoc@box\relax}%
  \Gscale@div\@tempb{\linewidth}{\wd\pandoc@box}%
  \ifdim\@tempb\p@<\@tempa\p@\let\@tempa\@tempb\fi% select the smaller of both
  \ifdim\@tempa\p@<\p@\scalebox{\@tempa}{\usebox\pandoc@box}%
  \else\usebox{\pandoc@box}%
  \fi%
}
% Set default figure placement to htbp
\def\fps@figure{htbp}
\makeatother
\usepackage{svg}
\ifLuaTeX
  \usepackage{luacolor}
  \usepackage[soul]{lua-ul}
\else
  \usepackage{soul}
\fi
\setlength{\emergencystretch}{3em} % prevent overfull lines
\providecommand{\tightlist}{%
  \setlength{\itemsep}{0pt}\setlength{\parskip}{0pt}}
\usepackage{bookmark}
\IfFileExists{xurl.sty}{\usepackage{xurl}}{} % add URL line breaks if available
\urlstyle{same}
\hypersetup{
  hidelinks,
  pdfcreator={LaTeX via pandoc}}

\author{}
\date{}

\begin{document}

\textbf{ВЫПУСКНАЯ КВАЛИФИКАЦИОННАЯ РАБОТА ПО ТЕМЕ\\
«РАЗРАБОТКА СИСТЕМЫ ДЛЯ ПРОКСИРОВАНИЯ И МОНИТОРИНГА СЕРВИСОВ-ЗАГЛУШЕК В
СРЕДЕ НАГРУЗОЧНОГО ТЕСТИРОВАНИЯ»}

{\def\LTcaptype{none} % do not increment counter
\begin{longtable}[]{@{}
  >{\raggedleft\arraybackslash}p{(\linewidth - 2\tabcolsep) * \real{0.5759}}
  >{\raggedleft\arraybackslash}p{(\linewidth - 2\tabcolsep) * \real{0.4241}}@{}}
\toprule\noalign{}
\begin{minipage}[b]{\linewidth}\raggedleft
Выполнил:
\end{minipage} & \begin{minipage}[b]{\linewidth}\raggedleft
Медведев Клим Николаевич
\end{minipage} \\
\midrule\noalign{}
\endhead
\bottomrule\noalign{}
\endlastfoot
Группа: & 221-329 \\
ТГ для связи: & @medvedevk \\
\end{longtable}
}

\textbf{\hfill\break
}

\section{СОДЕРЖАНИЕ}\label{ux441ux43eux434ux435ux440ux436ux430ux43dux438ux435}

\hyperref[ux43fux440ux435ux434ux43fux43eux43bux430ux433ux430ux435ux43cux430ux44f-ux441ux442ux440ux443ux43aux442ux443ux440ux430-ux440ux430ux431ux43eux442ux44b]{1
////// Предполагаемая структура работы //////
\hyperref[ux43fux440ux435ux434ux43fux43eux43bux430ux433ux430ux435ux43cux430ux44f-ux441ux442ux440ux443ux43aux442ux443ux440ux430-ux440ux430ux431ux43eux442ux44b]{4}}

\hyperref[ux43fux440ux435ux434ux43fux43eux43bux430ux433ux430ux435ux43cux43eux435-ux442ux435ux445ux43dux438ux447ux435ux441ux43aux43eux435-ux43eux43fux438ux441ux430ux43dux438ux435]{2
////// Предполагаемое техническое описание //////
\hyperref[ux43fux440ux435ux434ux43fux43eux43bux430ux433ux430ux435ux43cux43eux435-ux442ux435ux445ux43dux438ux447ux435ux441ux43aux43eux435-ux43eux43fux438ux441ux430ux43dux438ux435]{6}}

\hyperref[ux432ux432ux435ux434ux435ux43dux438ux435]{Введение
\hyperref[ux432ux432ux435ux434ux435ux43dux438ux435]{9}}

\hyperref[ux430ux43dux430ux43bux438ux437-ux43fux440ux435ux434ux43cux435ux442ux43dux43eux439-ux43eux431ux43bux430ux441ux442ux438-ux438-ux43eux431ux43eux441ux43dux43eux432ux430ux43dux438ux435-ux432ux44bux431ux43eux440ux430-ux442ux435ux445ux43dux43eux43bux43eux433ux438ux439]{3
АНАЛИЗ ПРЕДМЕТНОЙ ОБЛАСТИ И ОБОСНОВАНИЕ ВЫБОРА ТЕХНОЛОГИЙ
\hyperref[ux430ux43dux430ux43bux438ux437-ux43fux440ux435ux434ux43cux435ux442ux43dux43eux439-ux43eux431ux43bux430ux441ux442ux438-ux438-ux43eux431ux43eux441ux43dux43eux432ux430ux43dux438ux435-ux432ux44bux431ux43eux440ux430-ux442ux435ux445ux43dux43eux43bux43eux433ux438ux439]{12}}

\hyperref[ux438ux441ux441ux43bux435ux434ux43eux432ux430ux43dux438ux435-ux438ux43dux444ux440ux430ux441ux442ux440ux443ux43aux442ux443ux440ux44b-ux43dux430ux433ux440ux443ux437ux43eux447ux43dux43eux433ux43e-ux442ux435ux441ux442ux438ux440ux43eux432ux430ux43dux438ux44f-ux438-ux440ux43eux43bux438-ux438ux43cux438ux442ux430ux446ux438ux43eux43dux43dux44bux445-ux441ux435ux440ux432ux438ux441ux43eux432]{3.1
Исследование инфраструктуры нагрузочного тестирования и роли
имитационных сервисов
\hyperref[ux438ux441ux441ux43bux435ux434ux43eux432ux430ux43dux438ux435-ux438ux43dux444ux440ux430ux441ux442ux440ux443ux43aux442ux443ux440ux44b-ux43dux430ux433ux440ux443ux437ux43eux447ux43dux43eux433ux43e-ux442ux435ux441ux442ux438ux440ux43eux432ux430ux43dux438ux44f-ux438-ux440ux43eux43bux438-ux438ux43cux438ux442ux430ux446ux438ux43eux43dux43dux44bux445-ux441ux435ux440ux432ux438ux441ux43eux432]{12}}

\hyperref[ux430ux43dux430ux43bux438ux437-ux441ux443ux449ux435ux441ux442ux432ux443ux44eux449ux438ux445-ux43fux43eux434ux445ux43eux434ux43eux432-ux43a-ux443ux43fux440ux430ux432ux43bux435ux43dux438ux44e-ux441ux435ux440ux432ux438ux441ux430ux43cux438-ux437ux430ux433ux43bux443ux448ux43aux430ux43cux438-ux438-ux43fux440ux43eux431ux43bux435ux43cux430-ux438ux43cux43fux43eux440ux442ux43eux437ux430ux43cux435ux449ux435ux43dux438ux44f]{3.2
Анализ существующих подходов к управлению сервисами-заглушками и
проблема импортозамещения
\hyperref[ux430ux43dux430ux43bux438ux437-ux441ux443ux449ux435ux441ux442ux432ux443ux44eux449ux438ux445-ux43fux43eux434ux445ux43eux434ux43eux432-ux43a-ux443ux43fux440ux430ux432ux43bux435ux43dux438ux44e-ux441ux435ux440ux432ux438ux441ux430ux43cux438-ux437ux430ux433ux43bux443ux448ux43aux430ux43cux438-ux438-ux43fux440ux43eux431ux43bux435ux43cux430-ux438ux43cux43fux43eux440ux442ux43eux437ux430ux43cux435ux449ux435ux43dux438ux44f]{13}}

\hyperref[ux441ux440ux430ux432ux43dux438ux442ux435ux43bux44cux43dux44bux439-ux430ux43dux430ux43bux438ux437-ux438ux43dux441ux442ux440ux443ux43cux435ux43dux442ux430ux43bux44cux43dux44bux445-ux441ux440ux435ux434ux441ux442ux432-ux440ux430ux437ux440ux430ux431ux43eux442ux43aux438-ux438-ux43eux431ux43eux441ux43dux43eux432ux430ux43dux438ux435-ux442ux435ux445ux43dux43eux43bux43eux433ux438ux447ux435ux441ux43aux43eux433ux43e-ux441ux442ux435ux43aux430]{3.3
Сравнительный анализ инструментальных средств разработки и обоснование
технологического стека
\hyperref[ux441ux440ux430ux432ux43dux438ux442ux435ux43bux44cux43dux44bux439-ux430ux43dux430ux43bux438ux437-ux438ux43dux441ux442ux440ux443ux43cux435ux43dux442ux430ux43bux44cux43dux44bux445-ux441ux440ux435ux434ux441ux442ux432-ux440ux430ux437ux440ux430ux431ux43eux442ux43aux438-ux438-ux43eux431ux43eux441ux43dux43eux432ux430ux43dux438ux435-ux442ux435ux445ux43dux43eux43bux43eux433ux438ux447ux435ux441ux43aux43eux433ux43e-ux441ux442ux435ux43aux430]{15}}

\hyperref[ux43eux431ux43eux441ux43dux43eux432ux430ux43dux438ux435-ux432ux44bux431ux43eux440ux430-ux44fux437ux44bux43aux430-ux43fux440ux43eux433ux440ux430ux43cux43cux438ux440ux43eux432ux430ux43dux438ux44f-ux441ux435ux440ux432ux435ux440ux43dux43eux439-ux447ux430ux441ux442ux438]{3.3.1
Обоснование выбора языка программирования серверной части
\hyperref[ux43eux431ux43eux441ux43dux43eux432ux430ux43dux438ux435-ux432ux44bux431ux43eux440ux430-ux44fux437ux44bux43aux430-ux43fux440ux43eux433ux440ux430ux43cux43cux438ux440ux43eux432ux430ux43dux438ux44f-ux441ux435ux440ux432ux435ux440ux43dux43eux439-ux447ux430ux441ux442ux438]{16}}

\hyperref[ux441ux440ux430ux432ux43dux438ux442ux435ux43bux44cux43dux44bux439-ux430ux43dux430ux43bux438ux437-ux432ux435ux431-ux444ux440ux435ux439ux43cux432ux43eux440ux43aux43eux432]{3.3.2
Сравнительный анализ веб-фреймворков
\hyperref[ux441ux440ux430ux432ux43dux438ux442ux435ux43bux44cux43dux44bux439-ux430ux43dux430ux43bux438ux437-ux432ux435ux431-ux444ux440ux435ux439ux43cux432ux43eux440ux43aux43eux432]{18}}

\hyperref[ux43eux431ux43eux441ux43dux43eux432ux430ux43dux438ux435-ux432ux44bux431ux43eux440ux430-ux441ux438ux441ux442ux435ux43cux44b-ux443ux43fux440ux430ux432ux43bux435ux43dux438ux44f-ux434ux430ux43dux43dux44bux43cux438-ux441ux443ux431ux434]{3.3.3
Обоснование выбора системы управления данными (СУБД)
\hyperref[ux43eux431ux43eux441ux43dux43eux432ux430ux43dux438ux435-ux432ux44bux431ux43eux440ux430-ux441ux438ux441ux442ux435ux43cux44b-ux443ux43fux440ux430ux432ux43bux435ux43dux438ux44f-ux434ux430ux43dux43dux44bux43cux438-ux441ux443ux431ux434]{20}}

\hyperref[ux438ux442ux43eux433ux43eux432ux44bux439-ux432ux44bux431ux43eux440-ux441ux440ux435ux434ux441ux442ux432-ux440ux430ux437ux440ux430ux431ux43eux442ux43aux438]{3.3.4
Итоговый выбор средств разработки
\hyperref[ux438ux442ux43eux433ux43eux432ux44bux439-ux432ux44bux431ux43eux440-ux441ux440ux435ux434ux441ux442ux432-ux440ux430ux437ux440ux430ux431ux43eux442ux43aux438]{22}}

\hyperref[ux43eux431ux437ux43eux440-ux430ux43dux430ux43bux43eux433ux438ux447ux43dux44bux445-ux440ux435ux448ux435ux43dux438ux439-ux438-ux43eux431ux43eux441ux43dux43eux432ux430ux43dux438ux435-ux440ux430ux437ux440ux430ux431ux43eux442ux43aux438-ux441ux43eux431ux441ux442ux432ux435ux43dux43dux43eux433ux43e-ux43fux440ux43eux433ux440ux430ux43cux43cux43dux43eux433ux43e-ux43aux43eux43cux43fux43bux435ux43aux441ux430]{3.4
Обзор аналогичных решений и обоснование разработки собственного
программного комплекса
\hyperref[ux43eux431ux437ux43eux440-ux430ux43dux430ux43bux43eux433ux438ux447ux43dux44bux445-ux440ux435ux448ux435ux43dux438ux439-ux438-ux43eux431ux43eux441ux43dux43eux432ux430ux43dux438ux435-ux440ux430ux437ux440ux430ux431ux43eux442ux43aux438-ux441ux43eux431ux441ux442ux432ux435ux43dux43dux43eux433ux43e-ux43fux440ux43eux433ux440ux430ux43cux43cux43dux43eux433ux43e-ux43aux43eux43cux43fux43bux435ux43aux441ux430]{22}}

\hyperref[ux430ux43dux430ux43bux438ux437-ux440ux435ux448ux435ux43dux438ux44f-wiremock]{3.4.1
Анализ решения WireMock
\hyperref[ux430ux43dux430ux43bux438ux437-ux440ux435ux448ux435ux43dux438ux44f-wiremock]{23}}

\hyperref[ux430ux43dux430ux43bux438ux437-ux440ux435ux448ux435ux43dux438ux44f-mockserver]{3.4.2
Анализ решения MockServer
\hyperref[ux430ux43dux430ux43bux438ux437-ux440ux435ux448ux435ux43dux438ux44f-mockserver]{24}}

\hyperref[ux430ux43dux430ux43bux438ux437-ux440ux435ux448ux435ux43dux438ux44f-mountebank]{3.4.3
Анализ решения Mountebank
\hyperref[ux430ux43dux430ux43bux438ux437-ux440ux435ux448ux435ux43dux438ux44f-mountebank]{25}}

\hyperref[ux430ux43dux430ux43bux438ux437-ux440ux435ux448ux435ux43dux438ux44f-apache-tomcat]{3.4.4
Анализ решения Apache Tomcat
\hyperref[ux430ux43dux430ux43bux438ux437-ux440ux435ux448ux435ux43dux438ux44f-apache-tomcat]{26}}

\hyperref[ux441ux440ux430ux432ux43dux438ux442ux435ux43bux44cux43dux430ux44f-ux445ux430ux440ux430ux43aux442ux435ux440ux438ux441ux442ux438ux43aux430]{3.4.5
Сравнительная характеристика
\hyperref[ux441ux440ux430ux432ux43dux438ux442ux435ux43bux44cux43dux430ux44f-ux445ux430ux440ux430ux43aux442ux435ux440ux438ux441ux442ux438ux43aux430]{27}}

\hyperref[ux43eux431ux43eux441ux43dux43eux432ux430ux43dux438ux435-ux440ux430ux437ux440ux430ux431ux43eux442ux43aux438-ux441ux43eux431ux441ux442ux432ux435ux43dux43dux43eux433ux43e-ux43fux440ux43eux433ux440ux430ux43cux43cux43dux43eux433ux43e-ux43aux43eux43cux43fux43bux435ux43aux441ux430]{3.4.6
Обоснование разработки собственного программного комплекса
\hyperref[ux43eux431ux43eux441ux43dux43eux432ux430ux43dux438ux435-ux440ux430ux437ux440ux430ux431ux43eux442ux43aux438-ux441ux43eux431ux441ux442ux432ux435ux43dux43dux43eux433ux43e-ux43fux440ux43eux433ux440ux430ux43cux43cux43dux43eux433ux43e-ux43aux43eux43cux43fux43bux435ux43aux441ux430]{28}}

\hyperref[ux432ux44bux432ux43eux434-ux43fux43e-ux43fux43eux434ux440ux430ux437ux434ux435ux43bux443]{3.4.7
Вывод по подразделу
\hyperref[ux432ux44bux432ux43eux434-ux43fux43e-ux43fux43eux434ux440ux430ux437ux434ux435ux43bux443]{29}}

\hyperref[ux444ux43eux440ux43cux438ux440ux43eux432ux430ux43dux438ux435-ux444ux443ux43dux43aux446ux438ux43eux43dux430ux43bux44cux43dux44bux445-ux438-ux442ux435ux445ux43dux438ux447ux435ux441ux43aux438ux445-ux442ux440ux435ux431ux43eux432ux430ux43dux438ux439-ux43a-ux441ux438ux441ux442ux435ux43cux435]{3.5
Формирование функциональных и технических требований к системе
\hyperref[ux444ux43eux440ux43cux438ux440ux43eux432ux430ux43dux438ux435-ux444ux443ux43dux43aux446ux438ux43eux43dux430ux43bux44cux43dux44bux445-ux438-ux442ux435ux445ux43dux438ux447ux435ux441ux43aux438ux445-ux442ux440ux435ux431ux43eux432ux430ux43dux438ux439-ux43a-ux441ux438ux441ux442ux435ux43cux435]{30}}

\hyperref[ux442ux440ux435ux431ux43eux432ux430ux43dux438ux44f-ux43a-ux440ux435ux436ux438ux43cux430ux43c-ux444ux443ux43dux43aux446ux438ux43eux43dux438ux440ux43eux432ux430ux43dux438ux44f]{3.5.1
Требования к режимам функционирования
\hyperref[ux442ux440ux435ux431ux43eux432ux430ux43dux438ux44f-ux43a-ux440ux435ux436ux438ux43cux430ux43c-ux444ux443ux43dux43aux446ux438ux43eux43dux438ux440ux43eux432ux430ux43dux438ux44f]{30}}

\hyperref[ux444ux443ux43dux43aux446ux438ux43eux43dux430ux43bux44cux43dux44bux435-ux442ux440ux435ux431ux43eux432ux430ux43dux438ux44f]{3.5.2
Функциональные требования
\hyperref[ux444ux443ux43dux43aux446ux438ux43eux43dux430ux43bux44cux43dux44bux435-ux442ux440ux435ux431ux43eux432ux430ux43dux438ux44f]{31}}

\hyperref[ux43dux435ux444ux443ux43dux43aux446ux438ux43eux43dux430ux43bux44cux43dux44bux435-ux442ux435ux445ux43dux438ux447ux435ux441ux43aux438ux435-ux442ux440ux435ux431ux43eux432ux430ux43dux438ux44f]{3.5.3
Нефункциональные (технические) требования
\hyperref[ux43dux435ux444ux443ux43dux43aux446ux438ux43eux43dux430ux43bux44cux43dux44bux435-ux442ux435ux445ux43dux438ux447ux435ux441ux43aux438ux435-ux442ux440ux435ux431ux43eux432ux430ux43dux438ux44f]{32}}

\hyperref[ux442ux440ux435ux431ux43eux432ux430ux43dux438ux44f-ux43a-ux43fux440ux43eux433ux440ux430ux43cux43cux43dux43e-ux430ux43fux43fux430ux440ux430ux442ux43dux43eux43cux443-ux43eux431ux435ux441ux43fux435ux447ux435ux43dux438ux44e]{3.5.4
Требования к программно-аппаратному обеспечению
\hyperref[ux442ux440ux435ux431ux43eux432ux430ux43dux438ux44f-ux43a-ux43fux440ux43eux433ux440ux430ux43cux43cux43dux43e-ux430ux43fux43fux430ux440ux430ux442ux43dux43eux43cux443-ux43eux431ux435ux441ux43fux435ux447ux435ux43dux438ux44e]{33}}

\hyperref[ux442ux440ux435ux431ux43eux432ux430ux43dux438ux44f-ux43a-ux438ux43dux444ux43eux440ux43cux430ux446ux438ux43eux43dux43dux43eux439-ux431ux435ux437ux43eux43fux430ux441ux43dux43eux441ux442ux438]{3.5.5
Требования к информационной безопасности
\hyperref[ux442ux440ux435ux431ux43eux432ux430ux43dux438ux44f-ux43a-ux438ux43dux444ux43eux440ux43cux430ux446ux438ux43eux43dux43dux43eux439-ux431ux435ux437ux43eux43fux430ux441ux43dux43eux441ux442ux438]{33}}

\hyperref[ux432ux44bux432ux43eux434-ux43fux43e-ux440ux430ux437ux434ux435ux43bux443]{3.6
Вывод по разделу
\hyperref[ux432ux44bux432ux43eux434-ux43fux43e-ux440ux430ux437ux434ux435ux43bux443]{34}}

\hyperref[ux43fux440ux43eux435ux43aux442ux438ux440ux43eux432ux430ux43dux438ux435-ux438-ux442ux435ux445ux43dux438ux447ux435ux441ux43aux430ux44f-ux440ux435ux430ux43bux438ux437ux430ux446ux438ux44f-ux441ux438ux441ux442ux435ux43cux44b-ux443ux43fux440ux430ux432ux43bux435ux43dux438ux44f]{4
ПРОЕКТИРОВАНИЕ И ТЕХНИЧЕСКАЯ РЕАЛИЗАЦИЯ СИСТЕМЫ УПРАВЛЕНИЯ
\hyperref[ux43fux440ux43eux435ux43aux442ux438ux440ux43eux432ux430ux43dux438ux435-ux438-ux442ux435ux445ux43dux438ux447ux435ux441ux43aux430ux44f-ux440ux435ux430ux43bux438ux437ux430ux446ux438ux44f-ux441ux438ux441ux442ux435ux43cux44b-ux443ux43fux440ux430ux432ux43bux435ux43dux438ux44f]{35}}

\hyperref[ux440ux430ux437ux440ux430ux431ux43eux442ux43aux430-ux430ux440ux445ux438ux442ux435ux43aux442ux443ux440ux44b-ux440ux430ux441ux43fux440ux435ux434ux435ux43bux451ux43dux43dux43eux433ux43e-ux43fux440ux43eux433ux440ux430ux43cux43cux43dux43eux433ux43e-ux43aux43eux43cux43fux43bux435ux43aux441ux430]{4.1
Разработка архитектуры распределённого программного комплекса
\hyperref[ux440ux430ux437ux440ux430ux431ux43eux442ux43aux430-ux430ux440ux445ux438ux442ux435ux43aux442ux443ux440ux44b-ux440ux430ux441ux43fux440ux435ux434ux435ux43bux451ux43dux43dux43eux433ux43e-ux43fux440ux43eux433ux440ux430ux43cux43cux43dux43eux433ux43e-ux43aux43eux43cux43fux43bux435ux43aux441ux430]{35}}

\hyperref[ux432ux44bux431ux43eux440-ux430ux440ux445ux438ux442ux435ux43aux442ux443ux440ux43dux43eux433ux43e-ux441ux442ux438ux43bux44f]{4.1.1
Выбор архитектурного стиля
\hyperref[ux432ux44bux431ux43eux440-ux430ux440ux445ux438ux442ux435ux43aux442ux443ux440ux43dux43eux433ux43e-ux441ux442ux438ux43bux44f]{35}}

\hyperref[ux442ux43eux43fux43eux43bux43eux433ux438ux44f-ux440ux430ux437ux432ux451ux440ux442ux44bux432ux430ux43dux438ux44f-ux438-ux433ux438ux431ux43aux43eux441ux442ux44c-ux438ux43dux444ux440ux430ux441ux442ux440ux443ux43aux442ux443ux440ux44b]{4.1.2
Топология развёртывания и гибкость инфраструктуры
\hyperref[ux442ux43eux43fux43eux43bux43eux433ux438ux44f-ux440ux430ux437ux432ux451ux440ux442ux44bux432ux430ux43dux438ux44f-ux438-ux433ux438ux431ux43aux43eux441ux442ux44c-ux438ux43dux444ux440ux430ux441ux442ux440ux443ux43aux442ux443ux440ux44b]{38}}

\hyperref[ux43aux43eux43dux442ux435ux43aux441ux442ux43dux430ux44f-ux434ux438ux430ux433ux440ux430ux43cux43cux430-ux441ux438ux441ux442ux435ux43cux44b]{4.1.3
Контекстная диаграмма системы
\hyperref[ux43aux43eux43dux442ux435ux43aux441ux442ux43dux430ux44f-ux434ux438ux430ux433ux440ux430ux43cux43cux430-ux441ux438ux441ux442ux435ux43cux44b]{39}}

\hyperref[ux43aux43eux43dux442ux435ux439ux43dux435ux440ux43dux430ux44f-ux434ux438ux430ux433ux440ux430ux43cux43cux430]{4.1.4
Контейнерная диаграмма
\hyperref[ux43aux43eux43dux442ux435ux439ux43dux435ux440ux43dux430ux44f-ux434ux438ux430ux433ux440ux430ux43cux43cux430]{43}}

\hyperref[ux43aux43eux43cux43fux43eux43dux435ux43dux442ux43dux430ux44f-ux434ux438ux430ux433ux440ux430ux43cux43cux430-rest-api]{4.1.5
Компонентная диаграмма REST API
\hyperref[ux43aux43eux43cux43fux43eux43dux435ux43dux442ux43dux430ux44f-ux434ux438ux430ux433ux440ux430ux43cux43cux430-rest-api]{45}}

\hyperref[ux430ux440ux445ux438ux442ux435ux43aux442ux443ux440ux430-ux43cux43eux434ux443ux43bux44f-ux43dux430ux441ux442ux440ux43eux435ux43a-ux441ux438ux441ux442ux435ux43cux44b]{4.1.6
Архитектура модуля настроек системы
\hyperref[ux430ux440ux445ux438ux442ux435ux43aux442ux443ux440ux430-ux43cux43eux434ux443ux43bux44f-ux43dux430ux441ux442ux440ux43eux435ux43a-ux441ux438ux441ux442ux435ux43cux44b]{47}}

\hyperref[ux434ux438ux430ux433ux440ux430ux43cux43cux430-ux440ux430ux437ux432ux451ux440ux442ux44bux432ux430ux43dux438ux44f]{4.1.7
Диаграмма развёртывания
\hyperref[ux434ux438ux430ux433ux440ux430ux43cux43cux430-ux440ux430ux437ux432ux451ux440ux442ux44bux432ux430ux43dux438ux44f]{51}}

\hyperref[ux441ux446ux435ux43dux430ux440ux438ux439-ux432ux437ux430ux438ux43cux43eux434ux435ux439ux441ux442ux432ux438ux44f-ux43aux43eux43cux43fux43eux43dux435ux43dux442ux43eux432]{4.1.8
Сценарий взаимодействия компонентов
\hyperref[ux441ux446ux435ux43dux430ux440ux438ux439-ux432ux437ux430ux438ux43cux43eux434ux435ux439ux441ux442ux432ux438ux44f-ux43aux43eux43cux43fux43eux43dux435ux43dux442ux43eux432]{52}}

\hyperref[ux43fux440ux43eux435ux43aux442ux438ux440ux43eux432ux430ux43dux438ux435-ux441ux442ux440ux443ux43aux442ux443ux440-ux434ux430ux43dux43dux44bux445-ux434ux43bux44f-ux445ux440ux430ux43dux435ux43dux438ux44f-ux43aux43eux43dux444ux438ux433ux443ux440ux430ux446ux438ux439-ux438-ux441ux43eux441ux442ux43eux44fux43dux438ux439-ux432-redis]{4.2
Проектирование структур данных для хранения конфигураций и состояний в
Redis
\hyperref[ux43fux440ux43eux435ux43aux442ux438ux440ux43eux432ux430ux43dux438ux435-ux441ux442ux440ux443ux43aux442ux443ux440-ux434ux430ux43dux43dux44bux445-ux434ux43bux44f-ux445ux440ux430ux43dux435ux43dux438ux44f-ux43aux43eux43dux444ux438ux433ux443ux440ux430ux446ux438ux439-ux438-ux441ux43eux441ux442ux43eux44fux43dux438ux439-ux432-redis]{55}}

\hyperref[ux43eux431ux43eux441ux43dux43eux432ux430ux43dux438ux435-ux43cux43eux434ux435ux43bux438-ux434ux430ux43dux43dux44bux445]{4.2.1
Обоснование модели данных
\hyperref[ux43eux431ux43eux441ux43dux43eux432ux430ux43dux438ux435-ux43cux43eux434ux435ux43bux438-ux434ux430ux43dux43dux44bux445]{55}}

\hyperref[ux441ux442ux440ux443ux43aux442ux443ux440ux430-ux434ux430ux43dux43dux44bux445-ux437ux430ux433ux43bux443ux448ux43aux438]{4.2.2
Структура данных заглушки
\hyperref[ux441ux442ux440ux443ux43aux442ux443ux440ux430-ux434ux430ux43dux43dux44bux445-ux437ux430ux433ux43bux443ux448ux43aux438]{57}}

\hyperref[ux441ux442ux440ux443ux43aux442ux443ux440ux430-ux434ux430ux43dux43dux44bux445-ux445ux43eux441ux442ux430]{4.2.3
Структура данных хоста
\hyperref[ux441ux442ux440ux443ux43aux442ux443ux440ux430-ux434ux430ux43dux43dux44bux445-ux445ux43eux441ux442ux430]{59}}

\hyperref[ux441ux442ux440ux443ux43aux442ux443ux440ux430-ux434ux430ux43dux43dux44bux445-ux443ux447ux451ux442ux43dux43eux439-ux437ux430ux43fux438ux441ux438]{4.2.4
Структура данных учётной записи
\hyperref[ux441ux442ux440ux443ux43aux442ux443ux440ux430-ux434ux430ux43dux43dux44bux445-ux443ux447ux451ux442ux43dux43eux439-ux437ux430ux43fux438ux441ux438]{60}}

\hyperref[ux441ux442ux440ux443ux43aux442ux443ux440ux44b-ux434ux430ux43dux43dux44bux445-ux440ux435ux435ux441ux442ux440ux43eux432]{4.2.5
Структуры данных реестров
\hyperref[ux441ux442ux440ux443ux43aux442ux443ux440ux44b-ux434ux430ux43dux43dux44bux445-ux440ux435ux435ux441ux442ux440ux43eux432]{61}}

\hyperref[ux441ux442ux440ux443ux43aux442ux443ux440ux430-ux434ux430ux43dux43dux44bux445-ux441ux447ux451ux442ux447ux438ux43aux430-rate-limiter]{4.2.6
Структура данных счётчика Rate Limiter
\hyperref[ux441ux442ux440ux443ux43aux442ux443ux440ux430-ux434ux430ux43dux43dux44bux445-ux441ux447ux451ux442ux447ux438ux43aux430-rate-limiter]{63}}

\hyperref[ux433ux43bux43eux431ux430ux43bux44cux43dux44bux435-ux43dux430ux441ux442ux440ux43eux439ux43aux438-ux441ux438ux441ux442ux435ux43cux44b]{4.2.7
Глобальные настройки системы
\hyperref[ux433ux43bux43eux431ux430ux43bux44cux43dux44bux435-ux43dux430ux441ux442ux440ux43eux439ux43aux438-ux441ux438ux441ux442ux435ux43cux44b]{64}}

\hyperref[ux43fux435ux440ux441ux438ux441ux442ux435ux43dux442ux43dux43eux441ux442ux44c-ux434ux430ux43dux43dux44bux445-ux438-ux441ux442ux440ux430ux442ux435ux433ux438ux44f-ux432ux43eux441ux441ux442ux430ux43dux43eux432ux43bux435ux43dux438ux44f]{4.2.8
Персистентность данных и стратегия восстановления
\hyperref[ux43fux435ux440ux441ux438ux441ux442ux435ux43dux442ux43dux43eux441ux442ux44c-ux434ux430ux43dux43dux44bux445-ux438-ux441ux442ux440ux430ux442ux435ux433ux438ux44f-ux432ux43eux441ux441ux442ux430ux43dux43eux432ux43bux435ux43dux438ux44f]{65}}

\textbf{\hfill\break
}

\section{////// Предполагаемая структура работы
//////}\label{ux43fux440ux435ux434ux43fux43eux43bux430ux433ux430ux435ux43cux430ux44f-ux441ux442ux440ux443ux43aux442ux443ux440ux430-ux440ux430ux431ux43eux442ux44b}

\textbf{СОДЕРЖАНИЕ}

\textbf{ПЕРЕЧЕНЬ СОКРАЩЕНИЙ И ОБОЗНАЧЕНИЙ} (добавить: ASGI, AOF, RPS,
TTL)

\textbf{ВВЕДЕНИЕ}

\textbf{ГЛАВА 1. АНАЛИЗ ПРЕДМЕТНОЙ ОБЛАСТИ И ОБОСНОВАНИЕ ВЫБОРА
ТЕХНОЛОГИЙ}

1.1. Исследование инфраструктуры нагрузочного тестирования и роли
имитационных сервисов.

1.2. Анализ существующих подходов к управлению сервисами-заглушками и
проблема импортозамещения.

1.3. Сравнительный анализ и обоснование выбора технологического стека:
язык программирования, веб-фреймворк и система хранения данных

1.4. Обзор аналогичных решений и обоснование разработки собственного
программного комплекса.

1.5. Формирование функциональных и технических требований к системе.

\textbf{Выводы по первой главе.}

\textbf{ГЛАВА 2. ПРОЕКТИРОВАНИЕ И ТЕХНИЧЕСКАЯ РЕАЛИЗАЦИЯ СИСТЕМЫ
УПРАВЛЕНИЯ}

2.1. Разработка архитектуры распределенного программного комплекса.

2.2. Проектирование структур данных (Key-Value) для хранения
конфигураций и состояний процессов в In-Memory СУБД Redis.

2.3. Реализация асинхронных механизмов управления процессами и ввода
пользовательских параметров (JVM-аргументы).

2.4. Разработка модуля динамического проксирования запросов на базе
FastAPI и алгоритма Rate Limiting.

2.5. Проектирование и программная реализация пользовательского
интерфейса (Dashboard, управление кластером).

2.6. Реализация механизма экспорта метрик мониторинга и интеграция с
Prometheus.

\textbf{Выводы по второй главе.}

\textbf{\hl{ГЛАВА 3. ВНЕДРЕНИЕ, ТЕСТИРОВАНИЕ И}}
\textbf{\hl{ЭКСПЛУАТАЦИЯ РАЗРАБОТАННОГО ПРОГРАММНОГО КОМПЛЕКСА}}

3.1. Описание процесса развертывания системы в среде ОС Astra Linux
(настройка Uvicorn/Gunicorn и службы Redis).

3.2. Тестирование механизмов автоматического восстановления состояний
(персистентность Redis AOF) и перезапуска сервисов.

3.3. Экспериментальная оценка производительности прокси-модуля и
эффективности алгоритма Rate Limiting под нагрузкой.

3.4. Анализ работы системы в кластерном режиме и оценка алгоритмов
распределения нагрузки между узлами.

3.5. Рекомендации по дальнейшему развитию системы (шаблонизация
конфигураций, Self-healing).

\textbf{Выводы по третьей главе.}

\textbf{ЗАКЛЮЧЕНИЕ}

\textbf{СПИСОК ИСПОЛЬЗОВАННЫХ ИСТОЧНИКОВ}

\textbf{ПРИЛОЖЕНИЯ}

\begin{itemize}
\item
  Приложение А. Техническое задание.
\item
  Приложение Б. Программа и методика испытаний.
\item
  Приложение В. Листинги программного кода (реализация асинхронных
  эндпоинтов FastAPI, модуль работы с Redis).
\item
  Приложение Г. Протоколы тестирования и скриншоты интерфейса.
\end{itemize}

\section{////// Предполагаемое техническое описание
//////}\label{ux43fux440ux435ux434ux43fux43eux43bux430ux433ux430ux435ux43cux43eux435-ux442ux435ux445ux43dux438ux447ux435ux441ux43aux43eux435-ux43eux43fux438ux441ux430ux43dux438ux435}

\textbf{Тема проекта:} Разработка автоматизированной информационной
системы управления распределенной инфраструктурой имитационных сервисов
для нагрузочного тестирования.

\textbf{1. Назначение и концепция системы} Система предназначена для
централизованного управления жизненным циклом имитационных сервисов
(Java-приложений в форматах .jar и .war), используемых в качестве
«заглушек» внешних интеграций при проведении нагрузочного тестирования.
Программный комплекс позволяет автоматизировать процессы развертывания,
контроля состояния и маршрутизации трафика, обеспечивая
кросс-платформенную совместимость (в том числе с ОС Astra Linux) и
импортозамещение зарубежных решений.

\textbf{2. Гибкая архитектура и масштабируемость} Система поддерживает
два режима функционирования:

\begin{itemize}
\item
  \textbf{Локальный режим:} Все компоненты (управляющее приложение, СУБД
  Redis, имитационные сервисы) работают на одном хосте.
\item
  \textbf{Распределённый режим}: Управляющий узел взаимодействует с
  удалёнными хостами напрямую по протоколам SSH/SCP, без установки
  каких-либо компонентов системы на целевые серверы. Запуск, остановка и
  мониторинг процессов заглушек осуществляются через SSH-команды.
\end{itemize}

\textbf{3. Функциональные механизмы и алгоритмы} Реализация системы
предполагает использование следующих ключевых механизмов:

\begin{itemize}
\item
  \textbf{Управление процессами и транспорт артефактов:} Запуск
  имитационных сервисов на удалённых хостах инициируется напрямую через
  SSH-команды посредством библиотеки paramiko. Доставка исполняемых
  файлов (.jar) с Master-узла на Slave-узлы осуществляется по
  защищенному протоколу SSH (SCP), что гарантирует безопасность передачи
  данных.
\item
  \textbf{Динамическое управление портами:} При запуске нового сервиса
  система автоматически определяет свободный сетевой порт на целевом
  хосте. Информация о запущенном \hl{экземпляре приложения}
  (соответствие mock\_name -\textgreater{} ip:port) мгновенно
  регистрируется в Redis для настройки маршрутизации.
\item
  \textbf{Динамическая маршрутизация (Reverse Proxy):} Master-узел
  перенаправляет входящие запросы по URL-префиксу
  /\textless mock\_name\textgreater/ на соответствующий IP и порт
  сервиса, используя данные из Redis. Асинхронная модель ввода-вывода
  (Async I/O) FastAPI позволяет обрабатывать тысячи одновременных
  соединений с минимальными задержками.
\item
  \textbf{Алгоритм Rate Limiting (Fixed Window):} Для ограничения
  интенсивности запросов используется алгоритм «Счетчик с фиксированным
  окном» (Fixed Window Counter). При активации функции система
  инкрементирует атомарные счетчики в Redis с заданным временем жизни
  (TTL). При превышении порога (RPS) запросы блокируются с кодом 429 Too
  Many Requests.
\end{itemize}

\textbf{4. Работа с данными и сохранение состояний} Информационное
обеспечение системы базируется на высокопроизводительной In-Memory СУБД
\textbf{Redis}:

\begin{itemize}
\item
  \textbf{Хранение конфигураций:} Идентификаторы процессов (PID),
  параметры запуска, маппинг портов и метаданные сервисов хранятся в
  оперативной памяти (Key-Value), обеспечивая доступ менее чем за 1 мс.
\item
  \textbf{Персистентность состояний:} Для предотвращения потери данных
  используется механизм \textbf{AOF (Append Only File)}. Система
  фиксирует изменения на диске, что позволяет при перезагрузке
  управляющего приложения автоматически восстановить конфигурацию
  кластера.
\end{itemize}

\textbf{5. Мониторинг и метрики} Система реализует экспорт телеметрии
для сбора внешними средствами (Prometheus):

\begin{itemize}
\item
  \textbf{Эндпоинт:} /metrics.
\item
  \textbf{Метрики:} Текущий RPS по каждой заглушке, количество
  срабатываний Rate Limiter, статус доступности сервисов (UP/DOWN) и
  статус доступности целевых хостов.
\end{itemize}

\textbf{6. Пользовательский интерфейс} Интерфейс реализован с
использованием технологии \textbf{Server-Side Rendering (SSR)} на базе
шаблонизатора \textbf{Jinja2}. Это обеспечивает высокую скорость
отрисовки страниц, минимальные требования к клиенту и упрощает
архитектуру приложения (отсутствие отдельного SPA-фронтенда).

\begin{itemize}
\item
  \textbf{Главная панель (Dashboard):} Таблица со списком заглушек, их
  статусами, адресами и элементами управления.
\item
  \textbf{Управление кластером:} Форма добавления Slave-узлов и
  мониторинга их доступности.
\item
  \textbf{Редактор:} Интерфейс для загрузки .jar файлов, настройки
  шаблонов запуска (JVM-аргументы) и просмотра логов.
\end{itemize}

\section{Введение}\label{ux432ux432ux435ux434ux435ux43dux438ux435}

\textbf{Актуальность темы.} В современной практике обеспечения качества
сложных программных комплексов проведение нагрузочного тестирования
требует создания стабильной и управляемой среды. Одной из ключевых
проблем является необходимость имитации внешних систем, с которыми
взаимодействует объект тестирования. Использование имитационных сервисов
(«заглушек») позволяет изолировать тестируемую систему, однако
управление большим количеством таких сервисов, их развертывание на
удаленных серверах и контроль их состояния в ручном режиме существенно
замедляют процесс подготовки испытаний.

В условиях реализации политики импортозамещения в Российской Федерации
актуальность приобретает разработка отечественных инструментов
управления инфраструктурой, способных заменить зарубежные программные
продукты и обеспечить полноценное функционирование стендов нагрузочного
тестирования в среде защищенных операционных систем, таких как ОС Astra
Linux. Таким образом, создание системы, автоматизирующей жизненный цикл
имитационных сервисов и предоставляющей механизмы контроля трафика,
является своевременной и практически значимой задачей.

\hl{\textbf{Цель работы} -- проектирование и разработка
автоматизированной информационной системы для централизованного
управления и мониторинга состояния распределенной инфраструктуры
имитационных сервисов.}

\textbf{Задачи исследования:}

\begin{enumerate}
\def\labelenumi{\arabic{enumi}.}
\item
  Провести анализ предметной области и существующих подходов к
  управлению имитационными средами в процессах обеспечения качества ПО.
\item
  Спроектировать архитектуру программного комплекса, обеспечивающую
  гибкое масштабирование и работу в различных режимах.
\item
  Обосновать выбор технологического стека и программных средств
  реализации системы.
\item
  Разработать алгоритмы динамического проксирования трафика, механизмы
  контроля интенсивности запросов и программные интерфейсы для
  автоматизации управления.
\item
  Реализовать систему хранения конфигураций и контроля состояний
  исполняемых файлов.
\item
  Провести тестирование разработанного комплекса в среде ОС Astra Linux
  и оценить эффективность реализованных механизмов мониторинга и
  управления.
\end{enumerate}

\hl{\textbf{Объект исследования} -- инфраструктура запуска и
эксплуатации имитационных сервисов (заглушек) в среде нагрузочного
тестирования.}

\hl{\textbf{Предмет исследования} --} \hl{механизмы и алгоритмы
автоматизированного управления и контроля состояния имитационных
сервисов в среде нагрузочного тестирования.}

\textbf{Научная новизна} работы заключается в разработке \hl{алгоритма
централизованного управления} распределенной средой, поддерживающего
индивидуальную настройку ограничений интенсивности запросов (Rate
Limiting) для каждого сервиса, а также в реализации программного
интерфейса \hl{(Application Programming Interface,} \hl{API)} для
автоматического контроля состояния заглушек непосредственно из скриптов
нагрузочного тестирования.

\textbf{Практическая значимость} работы состоит в создании программного
инструмента, позволяющего отказаться от использования зарубежных систем
управления, упростить эксплуатацию инфраструктуры имитационных сервисов
и повысить уровень автоматизации процесса нагрузочного тестирования.

\textbf{Положения, выносимые на защиту:}

\begin{enumerate}
\def\labelenumi{\arabic{enumi}.}
\item
  Архитектура программного комплекса, поддерживающая гибкое
  масштабирование системы от локального запуска на одном хосте до
  формирования распределенного кластера.
\item
  Механизм динамического проксирования запросов, реализующий прозрачную
  маршрутизацию входящего трафика к целевым имитационным сервисам.
\item
  Алгоритм контроля интенсивности запросов (Rate Limiting),
  обеспечивающий стабильность работы имитационной среды при пиковых
  нагрузках.
\item
  Методика формирования и экспорта метрик производительности для
  непрерывного мониторинга состояния инфраструктуры заглушек.
\end{enumerate}

Структура работы. Выпускная квалификационная работа включает введение,
три главы основного текста, заключение, список использованных источников
и приложения.

Во введении обоснована актуальность исследования, поставлены цель и
задачи, определены объект и предмет работы, указаны научная новизна и
практическая значимость полученных результатов.

В первой главе проводится анализ предметной области, исследуются
существующие аналоги и методы управления имитационными средами. На
основе сравнительного анализа обосновывается выбор технологий и
формируются требования к системе.

Во второй главе описывается проектирование архитектуры системы,
структуры хранения данных и логики работы основных модулей. Приводятся
схемы взаимодействия компонентов и алгоритмы реализации функций
проксирования и управления кластером.

В третьей главе рассматриваются вопросы практической реализации,
настройки и эксплуатации системы в среде ОС Astra Linux. Приводятся
результаты тестирования, подтверждающие корректность работы
разработанных механизмов управления и мониторинга.

В заключении сформулированы основные выводы, подтверждающие решение
поставленных задач и достижение цели работы.

\section{АНАЛИЗ ПРЕДМЕТНОЙ ОБЛАСТИ И ОБОСНОВАНИЕ ВЫБОРА
ТЕХНОЛОГИЙ}\label{ux430ux43dux430ux43bux438ux437-ux43fux440ux435ux434ux43cux435ux442ux43dux43eux439-ux43eux431ux43bux430ux441ux442ux438-ux438-ux43eux431ux43eux441ux43dux43eux432ux430ux43dux438ux435-ux432ux44bux431ux43eux440ux430-ux442ux435ux445ux43dux43eux43bux43eux433ux438ux439}

\subsection{Исследование инфраструктуры нагрузочного тестирования и роли
имитационных
сервисов}\label{ux438ux441ux441ux43bux435ux434ux43eux432ux430ux43dux438ux435-ux438ux43dux444ux440ux430ux441ux442ux440ux443ux43aux442ux443ux440ux44b-ux43dux430ux433ux440ux443ux437ux43eux447ux43dux43eux433ux43e-ux442ux435ux441ux442ux438ux440ux43eux432ux430ux43dux438ux44f-ux438-ux440ux43eux43bux438-ux438ux43cux438ux442ux430ux446ux438ux43eux43dux43dux44bux445-ux441ux435ux440ux432ux438ux441ux43eux432}

В современной практике обеспечения качества сложных программных
комплексов нагрузочное тестирование занимает критически важное место.
Оно позволяет гарантировать стабильность и управляемость информационной
системы при высоких эксплуатационных нагрузках. Одной из ключевых
проблем при организации таких испытаний является необходимость
взаимодействия объекта тестирования с множеством внешних систем и
сервисов.

Использование реальных внешних интеграций в среде тестирования часто
оказывается невозможным или нецелесообразным по ряду причин:

\begin{itemize}
\item
  \textbf{Ограниченность ресурсов}: внешние системы могут иметь жесткие
  лимиты на количество запросов или требовать значительных финансовых
  затрат на эксплуатацию.
\item
  \textbf{Изоляция среды}: для получения достоверных результатов
  необходимо исключить влияние задержек и сбоев сторонних систем на
  показатели исследуемого объекта.
\item
  \textbf{Риски изменения данных}: проведение тестов на реальных
  интеграциях может привести к нежелательному изменению или повреждению
  информации в продуктовых базах данных.
\end{itemize}

Для решения этих проблем применяются \textbf{имитационные сервисы
(«заглушки»)} -- специализированное программное обеспечение, которое
воспроизводит поведение реальных компонентов, отвечая на запросы
тестируемой системы заранее определенным образом. \hl{В контексте данной
работы в качестве имитационных сервисов рассматриваются Java-приложения
в форматах \textbf{.jar} и \textbf{.war}.}

Роль имитационных сервисов в инфраструктуре нагрузочного тестирования
заключается в следующем:

\begin{enumerate}
\def\labelenumi{\arabic{enumi}.}
\item
  \textbf{Обеспечение автономности стенда}: заглушки позволяют полностью
  изолировать тестируемую систему от внешнего окружения, создавая
  предсказуемую среду.
\item
  \textbf{Эмуляция различных сценариев}: имитационные сервисы могут
  воспроизводить не только корректные ответы, но и ошибки (HTTP 5xx,
  тайм-ауты), что необходимо для проверки механизмов отказоустойчивости.
\item
  \textbf{Контроль интенсивности трафика}: использование механизмов
  ограничения запросов (Rate Limiting) позволяет моделировать поведение
  внешних систем при достижении ими предельных порогов
  производительности.
\end{enumerate}

Однако при увеличении масштаба информационной системы количество
необходимых заглушек возрастает, что порождает проблему управления
распределенной инфраструктурой. Процессы развертывания, контроля
состояния (активна/остановлена) и маршрутизации трафика к десяткам
имитационных сервисов на различных серверах в ручном режиме существенно
замедляют подготовку к испытаниям.

Таким образом, для эффективного функционирования стендов нагрузочного
тестирования необходима специализированная система. Она должна
централизованно управлять жизненным циклом заглушек, обеспечивать их
мониторинг и динамическое проксирование запросов, что особенно актуально
в условиях перехода на отечественные операционные системы, такие как
\textbf{ОС Astra Linux}.

\subsection{Анализ существующих подходов к управлению
сервисами-заглушками и проблема
импортозамещения}\label{ux430ux43dux430ux43bux438ux437-ux441ux443ux449ux435ux441ux442ux432ux443ux44eux449ux438ux445-ux43fux43eux434ux445ux43eux434ux43eux432-ux43a-ux443ux43fux440ux430ux432ux43bux435ux43dux438ux44e-ux441ux435ux440ux432ux438ux441ux430ux43cux438-ux437ux430ux433ux43bux443ux448ux43aux430ux43cux438-ux438-ux43fux440ux43eux431ux43bux435ux43cux430-ux438ux43cux43fux43eux440ux442ux43eux437ux430ux43cux435ux449ux435ux43dux438ux44f}

Традиционным подходом к управлению имитационными сервисами (заглушками),
разработанными на языке Java, является их развертывание в
специализированных контейнерах серверов приложений. В современной
практике разработки и тестирования сложился ряд стандартных решений,
среди которых наиболее распространенными являются:

\begin{enumerate}
\def\labelenumi{\arabic{enumi}.}
\item
  \textbf{Apache Tomcat} -- наиболее легкий и популярный
  сервлет-контейнер, используемый для запуска веб-приложений. Он
  обеспечивает базовую поддержку спецификаций Java Servlet и JavaServer
  Pages.
\item
  \textbf{WildFly (ранее JBoss Application Server)} -- более мощная и
  полнофункциональная платформа, поддерживающая широкий спектр
  стандартов Java EE (Jakarta EE). Данное решение предоставляет
  расширенные возможности управления ресурсами и кластеризации.
\item
  \textbf{Eclipse GlassFish} -- эталонная реализация платформы Java EE,
  предлагающая развитые инструменты администрирования через графическую
  консоль и командную строку.
\end{enumerate}

Анализ данных решений позволяет выделить общие тенденции и
характеристики классических серверов приложений:

\begin{itemize}
\item
  \textbf{Универсальность и полнофункциональность}: все перечисленные
  платформы спроектированы для запуска сложных корпоративных
  информационных систем. Однако для узкоспециализированной задачи --
  работы легковесных имитационных заглушек -- их функциональность
  является избыточной.
\item
  \textbf{Значительные накладные расходы}: общим недостатком является
  высокое потребление оперативной памяти и ресурсов центрального
  процессора самой платформой. При необходимости запуска десятков или
  сотен заглушек на одном или нескольких хостах суммарные затраты
  ресурсов на поддержание среды выполнения становятся критическими.
\item
  \textbf{Сложность автоматизации и оркестрации}: управление большим
  парком разрозненных экземпляров серверов приложений требует внедрения
  дополнительных тяжеловесных систем автоматизации. Штатные средства
  управления (консоли администрирования) ориентированы на ручное
  конфигурирование, что затрудняет их использование в динамических
  сценариях нагрузочного тестирования.
\item
  \textbf{Монолитная архитектура управления}: классические подходы
  подразумевают жесткую привязку приложения к конкретной конфигурации
  сервера, что снижает гибкость при необходимости быстрого изменения
  параметров запуска (например, JVM-аргументов) для отдельных сервисов.
\end{itemize}

На текущий момент в российской ИТ-индустрии наблюдается отчетливая
тенденция к отказу от использования зарубежных проприетарных и открытых
решений в пользу отечественного программного обеспечения или собственных
разработок. Это обусловлено политикой импортозамещения и необходимостью
обеспечения технологического суверенитета.

Большинство западных систем управления инфраструктурой и серверов
приложений (таких как продукты Oracle или Red Hat) сталкиваются с
проблемами поддержки, обновления и совместимости с российскими
операционными системами. В частности, эксплуатация имитационных стендов
в защищенной среде ОС Astra Linux требует инструментов, которые нативно
поддерживают механизмы управления этой операционной системы и не зависят
от зарубежных репозиториев и лицензий.

Таким образом, анализ существующих подходов показывает, что
использование традиционных серверов приложений типа Tomcat или WildFly
для управления распределенной сетью заглушек становится неэффективным.
Современные требования диктуют необходимость перехода к более легким,
гибким и отечественным решениям, способным обеспечить централизованное
управление жизненным циклом сервисов в рамках распределенных кластеров.

\subsection{Сравнительный анализ инструментальных средств разработки и
обоснование технологического
стека}\label{ux441ux440ux430ux432ux43dux438ux442ux435ux43bux44cux43dux44bux439-ux430ux43dux430ux43bux438ux437-ux438ux43dux441ux442ux440ux443ux43cux435ux43dux442ux430ux43bux44cux43dux44bux445-ux441ux440ux435ux434ux441ux442ux432-ux440ux430ux437ux440ux430ux431ux43eux442ux43aux438-ux438-ux43eux431ux43eux441ux43dux43eux432ux430ux43dux438ux435-ux442ux435ux445ux43dux43eux43bux43eux433ux438ux447ux435ux441ux43aux43eux433ux43e-ux441ux442ux435ux43aux430}

Выбор программных средств для реализации системы управления
имитационными сервисами является фундаментальным этапом проектирования.
От корректности этого выбора зависят эксплуатационные характеристики
будущей системы: предельная пропускная способность, латентность (время
отклика), масштабируемость и устойчивость к пиковым нагрузкам.

В ходе анализа предметной области и целей исследования было определено,
что разрабатываемая система должна обеспечивать бесперебойную
оркестрацию интенсивного потока запросов от нагрузочных станций. Это
накладывает жесткие ограничения на архитектуру: управляющий контур не
должен становиться «узким местом» (bottleneck) при маршрутизации
трафика. Требование минимизации накладных расходов исключает
использование технологий, склонных к блокировке вычислительных ресурсов
при ожидании ответа от внешних систем.

Для обоснования выбора проводится многокритериальный анализ по трем
направлениям: язык программирования, веб-фреймворк и система управления
базами данных.

\subsubsection{Обоснование выбора языка программирования серверной
части}\label{ux43eux431ux43eux441ux43dux43eux432ux430ux43dux438ux435-ux432ux44bux431ux43eux440ux430-ux44fux437ux44bux43aux430-ux43fux440ux43eux433ux440ux430ux43cux43cux438ux440ux43eux432ux430ux43dux438ux44f-ux441ux435ux440ux432ux435ux440ux43dux43eux439-ux447ux430ux441ux442ux438}

Для разработки backend-составляющей системы рассматривались три языка
программирования, наиболее востребованных в сфере высоконагруженных
систем (HighLoad) и системной автоматизации.

\begin{enumerate}
\def\labelenumi{\arabic{enumi}.}
\item
  \textbf{Java} -- объектно-ориентированный язык программирования.
  Является отраслевым стандартом для корпоративных систем и «родной»
  средой для запуска самих имитационных сервисов (.jar).
\end{enumerate}

\begin{itemize}
\item
  \textbf{Преимущества:}

  \begin{itemize}
  \item
    Строгая статическая типизация, снижающая количество ошибок на этапе
    компиляции.
  \item
    Развитая экосистема библиотек и инструментов профилирования.
  \item
    Поддержка полноценной многопоточности (multithreading).
  \end{itemize}
\item
  \textbf{Недостатки:}

  \begin{itemize}
  \item
    Высокая ресурсоемкость виртуальной машины (JVM), требующая
    значительного объема оперативной памяти.
  \item
    Длительное время «холодного старта», что критично для динамических
    сред тестирования.
  \item
    Громоздкий синтаксис для простых задач взаимодействия с процессами
    операционной системы.
  \end{itemize}
\end{itemize}

\begin{enumerate}
\def\labelenumi{\arabic{enumi}.}
\setcounter{enumi}{1}
\item
  \textbf{Go (Golang) --} компилируемый многопоточный язык
  программирования, разработанный Google, ориентированный на создание
  эффективных сетевых сервисов.
\end{enumerate}

\begin{itemize}
\item
  \textbf{Преимущества:}

  \begin{itemize}
  \item
    Высочайшая производительность, сопоставимая с C++, благодаря
    компиляции в нативный код.
  \item
    Эффективная модель конкурентности (goroutines), позволяющая
    обрабатывать тысячи соединений с минимальными накладными расходами.
  \item
    Отсутствие внешних зависимостей (компиляция в единый бинарный файл).
  \end{itemize}
\item
  \textbf{Недостатки:}

  \begin{itemize}
  \item
    Отсутствие гибкости при работе с динамическими структурами данных
    (JSON произвольной структуры) из-за строгой типизации.
  \item
    Специфическая обработка ошибок и отсутствие исключений, что
    увеличивает объем шаблонного кода (boilerplate).
  \item
    Более высокий порог входа по сравнению с интерпретируемыми языками.
  \end{itemize}
\end{itemize}

\begin{enumerate}
\def\labelenumi{\arabic{enumi}.}
\setcounter{enumi}{2}
\item
  \textbf{Python --} высокоуровневый интерпретируемый язык
  программирования общего назначения с динамической строгой типизацией.
\end{enumerate}

\begin{itemize}
\item
  \textbf{Преимущества:}

  \begin{itemize}
  \item
    Наличие мощных инструментов асинхронного программирования
    (библиотека asyncio), позволяющих эффективно обрабатывать задачи,
    ограниченные скоростью \hl{ввода-вывода (I/O-bound),} такие как
    сетевые запросы или ожидание ответа от базы данных.
  \item
    Глубокая интеграция с Linux: модуль subprocess предоставляет удобный
    интерфейс для управления процессами, сигналами и потоками
    ввода-вывода.
  \item
    Высокая скорость разработки благодаря лаконичному синтаксису.
  \end{itemize}
\item
  \textbf{Недостатки:}

  \begin{itemize}
  \item
    Производительность интерпретатора ниже, чем у компилируемых языков.
  \item
    Наличие Global Interpreter Lock (GIL), ограничивающего выполнение
    потоков на одном ядре CPU (нивелируется использованием асинхронного
    подхода).
  \end{itemize}
\end{itemize}

\textbf{Вывод:} В качестве основного языка выбран \textbf{Python}.
Возможность использования асинхронного программирования в сочетании с
нативными средствами управления процессами Linux позволяет обеспечить
необходимую производительность, сохраняя гибкость и скорость разработки.

\subsubsection{Сравнительный анализ
веб-фреймворков}\label{ux441ux440ux430ux432ux43dux438ux442ux435ux43bux44cux43dux44bux439-ux430ux43dux430ux43bux438ux437-ux432ux435ux431-ux444ux440ux435ux439ux43cux432ux43eux440ux43aux43eux432}

Учитывая необходимость обеспечения высокой пропускной способности
системы, выбор архитектуры веб-фреймворка является критическим.
Рассматривались три кандидата: Django, Flask и FastAPI.

\begin{enumerate}
\def\labelenumi{\arabic{enumi}.}
\item
  \textbf{Django --} полнофункциональный синхронный фреймворк уровня
  Enterprise, следующий принципу «Batteries included» (все включено).
\end{enumerate}

\begin{itemize}
\item
  \textbf{Преимущества:}

  \begin{itemize}
  \item
    Наличие встроенной панели администрирования, ORM (Object-Relational
    Mapping) и системы аутентификации «из коробки».
  \item
    Развитая экосистема плагинов и высокий уровень безопасности.
  \item
    Строгая структура проекта, облегчающая поддержку крупных монолитных
    приложений.
  \end{itemize}
\item
  \textbf{Недостатки:}

  \begin{itemize}
  \item
    Монолитная архитектура и избыточный функционал создают значительные
    накладные расходы (overhead) на каждый запрос.
  \item
    Синхронная модель обработки запросов плохо подходит для
    высоконагруженных микросервисов, так как блокирует потоки при
    ожидании ввода-вывода.
  \item
    Сложность отключения неиспользуемых компонентов для облегчения
    системы.
  \end{itemize}
\end{itemize}

\begin{enumerate}
\def\labelenumi{\arabic{enumi}.}
\setcounter{enumi}{1}
\item
  \textbf{Flask --} легковесный микрофреймворк. Классическое решение,
  основанное на спецификации \hl{WSGI (Web Server Gateway Interface).}
\end{enumerate}

\begin{itemize}
\item
  \textbf{Преимущества:}

  \begin{itemize}
  \item
    Минимализм и модульность: разработчик подключает только необходимые
    библиотеки.
  \item
    Низкий порог входа и простота создания простых сервисов.
  \item
    Гибкость в выборе архитектурных паттернов.
  \end{itemize}
\item
  \textbf{Недостатки:}

  \begin{itemize}
  \item
    Использование синхронной (блокирующей) модели ввода-вывода. При
    высокой нагрузке это приводит к исчерпанию пула потоков сервера и
    росту задержек.
  \item
    Отсутствие встроенных механизмов валидации данных (требует
    подключения сторонних расширений).
  \item
    Снижение производительности при большом количестве одновременных
    соединений (проблема C10k).
  \end{itemize}
\end{itemize}

\begin{enumerate}
\def\labelenumi{\arabic{enumi}.}
\setcounter{enumi}{2}
\item
  \textbf{FastAPI} -- современный асинхронный фреймворк. Инструмент,
  построенный на спецификации ASGI (Asynchronous Server Gateway
  Interface).
\end{enumerate}

\begin{itemize}
\item
  \textbf{Преимущества:}

  \begin{itemize}
  \item
    Полная поддержка асинхронности (async/await) и событийного цикла,
    что позволяет обрабатывать тысячи запросов в секунду на одном
    процессе.
  \item
    Высокая производительность благодаря использованию Starlette и
    uvloop, сопоставимая с Go и Node.js.
  \item
    Встроенная быстрая валидация данных и сериализация на базе Pydantic.
  \end{itemize}
\item
  \textbf{Недостатки:}

  \begin{itemize}
  \item
    Относительно молодая экосистема по сравнению с Django и Flask.
  \item
    Требует понимания принципов асинхронного программирования.
  \end{itemize}
\end{itemize}

\textbf{Вывод:} Выбран фреймворк \textbf{FastAPI}. Это единственное
решение в экосистеме Python, способное обеспечить стабильную обработку
интенсивного трафика с минимальными задержками благодаря асинхронной
неблокирующей архитектуре.

\subsubsection{Обоснование выбора системы управления данными
(СУБД)}\label{ux43eux431ux43eux441ux43dux43eux432ux430ux43dux438ux435-ux432ux44bux431ux43eux440ux430-ux441ux438ux441ux442ux435ux43cux44b-ux443ux43fux440ux430ux432ux43bux435ux43dux438ux44f-ux434ux430ux43dux43dux44bux43cux438-ux441ux443ux431ux434}

Для обеспечения высокой отзывчивости системы подсистема хранения данных
должна обладать минимальным временем отклика на чтение конфигураций и
запись статусов. Сравнивались конкретные распространенные решения.

\begin{enumerate}
\def\labelenumi{\arabic{enumi}.}
\item
  \textbf{PostgreSQL -- объектно-реляционная СУБД.} Мощная система баз
  данных с открытым исходным кодом.
\end{enumerate}

\begin{itemize}
\item
  \textbf{Преимущества:}

  \begin{itemize}
  \item
    Гарантия целостности данных (ACID) и надежность хранения.
  \item
    Мощный язык запросов SQL и поддержка сложных выборок (JOIN).
  \item
    Широкие возможности расширяемости (типы данных, индексы).
  \end{itemize}
\item
  \textbf{Недостатки:}

  \begin{itemize}
  \item
    Накладные расходы на дисковые операции ввода-вывода и журналирование
    (WAL) замедляют отклик.
  \item
    Избыточность функционала для хранения простых пар «ключ-значение»
    (конфигурации заглушек).
  \item
    Потребность в строгом проектировании схемы данных (Schema
    migrations).
  \end{itemize}
\end{itemize}

\begin{enumerate}
\def\labelenumi{\arabic{enumi}.}
\setcounter{enumi}{1}
\item
  \textbf{MySQL -- реляционная СУБД.} Популярная база данных, часто
  используемая в веб-приложениях.
\end{enumerate}

\begin{itemize}
\item
  \textbf{Преимущества:}

  \begin{itemize}
  \item
    Высокая скорость операций чтения в простых сценариях.
  \item
    Широкая распространенность и простота администрирования.
  \item
    Поддержка различных движков хранения (Storage Engines).
  \end{itemize}
\item
  \textbf{Недостатки:}

  \begin{itemize}
  \item
    Как и любая реляционная СУБД, зависит от скорости дисковой
    подсистемы.
  \item
    Блокировки таблиц или строк могут снижать конкурентность записи при
    высокой нагрузке.
  \item
    Жесткость схемы данных.
  \end{itemize}
\end{itemize}

\begin{enumerate}
\def\labelenumi{\arabic{enumi}.}
\setcounter{enumi}{2}
\item
  \textbf{MongoDB -- документоориентированная NoSQL СУБД}. Система,
  хранящая данные в формате BSON (бинарный JSON).
\end{enumerate}

\begin{itemize}
\item
  \textbf{Преимущества:}

  \begin{itemize}
  \item
    Гибкая схема данных (Schemaless), идеально подходящая для хранения
    JSON-конфигураций.
  \item
    Горизонтальная масштабируемость (шардинг) из коробки.
  \end{itemize}
\item
  \textbf{Недостатки:}

  \begin{itemize}
  \item
    Потребляет значительно больше памяти и дискового пространства по
    сравнению с аналогами.
  \item
    В стандартной конфигурации уступает In-Memory решениям по скорости
    отклика.
  \item
    Слабые гарантии транзакционности в старых версиях.
  \end{itemize}
\end{itemize}

\begin{enumerate}
\def\labelenumi{\arabic{enumi}.}
\setcounter{enumi}{3}
\item
  \textbf{Redis -- In-Memory хранилище структур данных}.
  Высокопроизводительная СУБД, работающая полностью в оперативной
  памяти.
\end{enumerate}

\begin{itemize}
\item
  \textbf{Преимущества:}

  \begin{itemize}
  \item
    Экстремально низкая латентность (время доступа менее 1 мс) благодаря
    отсутствию обращений к диску при чтении.
  \item
    Простая модель данных «Ключ-Значение», полностью соответствующая
    задаче хранения состояний процессов.
  \item
    Нативная поддержка механизма Time-To-Live (TTL) для временных
    данных.
  \end{itemize}
\item
  \textbf{Недостатки:}

  \begin{itemize}
  \item
    Зависимость объема хранимых данных от объема оперативной памяти.
  \item
    Риск потери данных при аварийном отключении питания (решается
    настройкой персистентности).
  \end{itemize}
\end{itemize}

\textbf{Вывод:} Выбрана СУБД \textbf{Redis}. Использование In-Memory
архитектуры обеспечивает максимальное быстродействие. Проблема
сохранности данных решается включением режима \textbf{AOF (Append Only
File)}, который асинхронно сохраняет операции на диск, гарантируя
восстановление конфигураций после перезагрузки без ущерба для
производительности.

\subsubsection{Итоговый выбор средств
разработки}\label{ux438ux442ux43eux433ux43eux432ux44bux439-ux432ux44bux431ux43eux440-ux441ux440ux435ux434ux441ux442ux432-ux440ux430ux437ux440ux430ux431ux43eux442ux43aux438}

На основании проведенного анализа и выявленных требований к
быстродействию и надежности, утвержден следующий стек технологий для
реализации выпускной квалификационной работы:

\begin{enumerate}
\def\labelenumi{\arabic{enumi}.}
\item
  \textbf{Язык программирования:} Python (с использованием asyncio).
\item
  \textbf{Веб-фреймворк:} FastAPI (сервер приложений Uvicorn).
\item
  \textbf{СУБД:} Redis (режим персистентности AOF).
\item
  \textbf{Целевая платформа:} ОС Astra Linux Special Edition.
\end{enumerate}

Данный выбор обеспечивает оптимальный баланс между производительностью,
скоростью разработки и надежностью функционирования системы в условиях
высокой нагрузки.

\subsection{Обзор аналогичных решений и обоснование разработки
собственного программного
комплекса}\label{ux43eux431ux437ux43eux440-ux430ux43dux430ux43bux43eux433ux438ux447ux43dux44bux445-ux440ux435ux448ux435ux43dux438ux439-ux438-ux43eux431ux43eux441ux43dux43eux432ux430ux43dux438ux435-ux440ux430ux437ux440ux430ux431ux43eux442ux43aux438-ux441ux43eux431ux441ux442ux432ux435ux43dux43dux43eux433ux43e-ux43fux440ux43eux433ux440ux430ux43cux43cux43dux43eux433ux43e-ux43aux43eux43cux43fux43bux435ux43aux441ux430}

В рамках анализа предметной области необходимо рассмотреть существующие
на рынке инструменты, применяемые для создания, эксплуатации и
управления имитационными сервисами (заглушками). Поскольку
разрабатываемая система предназначена для управления инфраструктурой
заглушек в процессе нагрузочного тестирования, к аналогам следует
отнести программные продукты, обеспечивающие виртуализацию сервисов
(Service Virtualization), мокирование (Mocking), а также серверы
приложений, традиционно используемые для хостинга Java-артефактов.

Сравнительный анализ проводится по критериям, критически важным для
обеспечения стабильности нагрузочного стенда и автоматизации процессов
тестирования:

\begin{enumerate}
\def\labelenumi{\arabic{enumi}.}
\item
  \textbf{Архитектура и потребление ресурсов:} способность работать в
  условиях ограниченных аппаратных мощностей при запуске десятков
  экземпляров заглушек.
\item
  \textbf{Возможности управления процессами (Orchestration):} наличие
  инструментов для удаленного запуска, остановки и обновления версий
  приложений-заглушек.
\item
  \textbf{Функциональность управления трафиком:} наличие встроенных
  механизмов ограничения нагрузки (Rate Limiting) и эмуляции сетевых
  задержек.
\item
  \textbf{Совместимость и импортозамещение:} возможность автономной
  работы в среде ОС Astra Linux без зависимости от зарубежных облачных
  сервисов и проприетарных лицензий.
\end{enumerate}

\subsubsection{Анализ решения
WireMock}\label{ux430ux43dux430ux43bux438ux437-ux440ux435ux448ux435ux43dux438ux44f-wiremock}

\textbf{WireMock} -- один из наиболее распространенных инструментов с
открытым исходным кодом для виртуализации HTTP-сервисов. Продукт
позволяет создавать заглушки, настраивать правила сопоставления запросов
(request matching) и формировать динамические ответы.

\textbf{Функциональные возможности:}

Согласно официальной документации {[}1{]}, WireMock поддерживает гибкую
настройку ответов через JSON API, имитацию задержек (network latency) и
ошибок (fault injection). Он может работать в двух режимах: как
библиотека, встраиваемая в код автотестов, или как отдельный процесс
(Standalone).

\textbf{Недостатки в контексте поставленной задачи:}

\begin{enumerate}
\def\labelenumi{\arabic{enumi}.}
\item
  Ресурсоемкость: WireMock разработан на Java и выполняется в
  виртуальной машине JVM. Запуск отдельного экземпляра WireMock для
  каждого имитируемого микросервиса создает значительные накладные
  расходы на оперативную память (RAM) и процессорное время (CPU). В
  условиях нагрузочного тестирования, когда требуется поднять 50--100
  заглушек, суммарные затраты ресурсов на поддержание работы самих
  заглушек становятся неприемлемыми.
\item
  Отсутствие встроенной оркестрации: В бесплатной версии (Open Source)
  WireMock отсутствует централизованная панель управления распределенным
  кластером. Инструмент фокусируется на логике ответа, а не на
  эксплуатации сервера.
\item
  Ограничения Rate Limiting: Базовый функционал позволяет имитировать
  задержки, однако полноценный алгоритм ограничения количества запросов
  в секунду (RPS) с отбрасыванием лишнего трафика требует написания
  кастомных расширений (extensions) на Java {[}1{]}.
\end{enumerate}

\subsubsection{Анализ решения
MockServer}\label{ux430ux43dux430ux43bux438ux437-ux440ux435ux448ux435ux43dux438ux44f-mockserver}

\textbf{MockServer} -- популярное решение для мокирования систем,
взаимодействующих по протоколам HTTP или HTTPS. Продукт часто
используется в контейнеризированных средах (Docker, Kubernetes) для
интеграционного тестирования.

\textbf{Функциональные возможности:}

MockServer предоставляет мощный API для создания «ожиданий»
(Expectations) и верификации запросов. Согласно документации {[}2{]}, он
поддерживает проксирование трафика (Port Forwarding), запись проходящих
запросов для последующего анализа и генерацию ответов на основе шаблонов
(Mustache, Velocity).

\textbf{Недостатки в контексте поставленной задачи:}

\begin{enumerate}
\def\labelenumi{\arabic{enumi}.}
\item
  Монолитная архитектура управления: MockServer хранит состояния
  «ожиданий» в памяти конкретного экземпляра приложения. Для создания
  распределенной системы требуется сложная настройка кластеризации, что
  утяжеляет инфраструктуру тестового стенда.
\item
  Сложность динамического управления: MockServer ориентирован на
  программное управление через API во время прогона тестов.
  Использование его как постоянно действующего инфраструктурного
  компонента требует разработки дополнительной обвязки для мониторинга
  его состояния и автоматического перезапуска (Self-healing) {[}2{]}.
\item
  Зависимость от Java: Аналогично WireMock, продукт написан на Java, что
  влечет проблемы с «холодным стартом» и высоким потреблением памяти
  (Heap Size), что снижает плотность размещения заглушек на сервере.
\end{enumerate}

\subsubsection{Анализ решения
Mountebank}\label{ux430ux43dux430ux43bux438ux437-ux440ux435ux448ux435ux43dux438ux44f-mountebank}

\textbf{Mountebank} -- инструмент с открытым исходным кодом,
разработанный на платформе Node.js. Его ключевой особенностью является
мультипротокольность: помимо HTTP/HTTPS, он поддерживает TCP, SMTP и
другие протоколы через систему плагинов.

\textbf{Функциональные возможности:}

Mountebank использует концепцию «самозванцев» (imposters) -- легковесных
сервисов, слушающих определенные порты. Документация {[}3{]} указывает
на возможность использования JavaScript-инъекций для реализации сложной
логики обработки ответов.

\textbf{Недостатки в контексте поставленной задачи:}

\begin{enumerate}
\def\labelenumi{\arabic{enumi}.}
\item
  Проблемы безопасности: Возможность исполнения произвольного JS-кода
  (injection) в теле запроса создает уязвимости в контуре нагрузочного
  тестирования.
\item
  Отсутствие управления артефактами: Mountebank позволяет создать логику
  ответа (скрипт), но он не предназначен для запуска и управления
  бинарными файлами (заглушками в виде .jar), предоставленными
  разработчиками.
\item
  Производительность в однопоточном режиме: В стандартной конфигурации
  Mountebank работает в одном потоке (Event Loop). При высокой нагрузке
  (тысячи RPS), характерной для нагрузочного тестирования, сложные
  вычисления внутри одной заглушки могут заблокировать обработку
  запросов других заглушек, работающих в том же экземпляре Mountebank
  {[}3{]}.
\end{enumerate}

\subsubsection{Анализ решения Apache
Tomcat}\label{ux430ux43dux430ux43bux438ux437-ux440ux435ux448ux435ux43dux438ux44f-apache-tomcat}

\textbf{Apache Tomcat} -- свободно распространяемый контейнер сервлетов,
реализующий спецификации Jakarta Servlet и Jakarta Pages. Это
классическое решение для хостинга Java-приложений, широко применяемое в
промышленной эксплуатации.

\textbf{Функциональные возможности:}

Согласно официальной документации {[}4{]}, Tomcat предоставляет развитые
средства управления жизненным циклом веб-приложений. Компонент Manager
App позволяет развертывать (deploy), останавливать и перезапускать
приложения без остановки сервера. Поддерживается кластеризация для
репликации сессий {[}5{]}. Для защиты от перегрузок предусмотрен фильтр
RateLimitFilter, позволяющий ограничивать количество запросов с одного
IP-адреса {[}6{]}.

\textbf{Недостатки в контексте поставленной задачи:}

\begin{enumerate}
\def\labelenumi{\arabic{enumi}.}
\item
  Избыточность архитектуры (Overhead): Tomcat является полноценным
  сервлет-контейнером. Использование его для запуска примитивных
  заглушек нерационально с точки зрения потребления ресурсов. Запуск
  отдельного экземпляра Tomcat для каждой заглушки потребляет сотни
  мегабайт оперативной памяти только на нужды самого сервера, что
  критично при дефиците ресурсов.
\item
  Требования к формату артефактов: Tomcat работает с WAR-архивами.
  Однако современные микросервисы (и заглушки для них) чаще всего
  поставляются в виде Fat JAR (со встроенным сервером, например, Netty
  или Undertow). Для запуска таких заглушек в Tomcat требуется их
  пересборка и модификация кода, что усложняет процесс.
\item
  Статичность конфигурации ограничений: Механизмы Rate Limiting в Tomcat
  конфигурируются декларативно через XML-файлы при старте. В задачах
  нагрузочного тестирования требуется \emph{динамическое} изменение
  параметров дросселирования трафика «на лету» (runtime) без
  перезагрузки сервера, что в Tomcat не реализовано «из коробки».
\end{enumerate}

\subsubsection{Сравнительная
характеристика}\label{ux441ux440ux430ux432ux43dux438ux442ux435ux43bux44cux43dux430ux44f-ux445ux430ux440ux430ux43aux442ux435ux440ux438ux441ux442ux438ux43aux430}

Для наглядного сравнения рассмотренных решений с проектируемой системой
составлена таблица 1.

Таблица 1 -- Сравнительный анализ средств управления имитационными
сервисами

{\def\LTcaptype{none} % do not increment counter
\begin{longtable}[]{@{}
  >{\raggedright\arraybackslash}p{(\linewidth - 10\tabcolsep) * \real{0.2094}}
  >{\raggedright\arraybackslash}p{(\linewidth - 10\tabcolsep) * \real{0.1509}}
  >{\raggedright\arraybackslash}p{(\linewidth - 10\tabcolsep) * \real{0.1509}}
  >{\raggedright\arraybackslash}p{(\linewidth - 10\tabcolsep) * \real{0.1485}}
  >{\raggedright\arraybackslash}p{(\linewidth - 10\tabcolsep) * \real{0.1296}}
  >{\raggedright\arraybackslash}p{(\linewidth - 10\tabcolsep) * \real{0.2107}}@{}}
\toprule\noalign{}
\begin{minipage}[b]{\linewidth}\centering
\textbf{Критерий сравнения}
\end{minipage} & \begin{minipage}[b]{\linewidth}\centering
\textbf{WireMock}
\end{minipage} & \begin{minipage}[b]{\linewidth}\centering
\textbf{MockServer}
\end{minipage} & \begin{minipage}[b]{\linewidth}\centering
\textbf{Mountebank}
\end{minipage} & \begin{minipage}[b]{\linewidth}\centering
\textbf{Apache Tomcat}
\end{minipage} & \begin{minipage}[b]{\linewidth}\centering
\textbf{Разрабатываемая система}
\end{minipage} \\
\midrule\noalign{}
\endhead
\bottomrule\noalign{}
\endlastfoot
Технологическая платформа & Java (JVM) & Java (JVM) & Node.js & Java
(JVM) & Python (FastAPI) + Redis \\
Тип архитектуры & Монолитный процесс & Монолитный процесс & Процесс
Node.js & Сервлет-контейнер & Агент управления процессами
(Master-Agent) \\
Потребление ресурсов (RAM/CPU) & Высокое & Высокое & Среднее & Очень
высокое & Низкое (асинхронная модель) \\
Управление запуском .jar файлов & Нет & Нет & Нет & Нет (только .war) &
Да (прямая поддержка) \\
Динамический Rate Limiting & Требует расширений & Ограничено &
Скриптовой (JS) & Статичный (XML) & Да (Управляемый через API) \\
Централизованное управление & Нет (в OSS версии) & Сложное (Docker/K8s)
& Нет & Сложное (Cluster) & Встроено (Dashboard) \\
\end{longtable}
}

\emph{Источники данных для таблицы:} \emph{\hl{{[}1{]}, {[}2{]},
{[}3{]}, {[}4{]}.}}

\subsubsection{Обоснование разработки собственного программного
комплекса}\label{ux43eux431ux43eux441ux43dux43eux432ux430ux43dux438ux435-ux440ux430ux437ux440ux430ux431ux43eux442ux43aux438-ux441ux43eux431ux441ux442ux432ux435ux43dux43dux43eux433ux43e-ux43fux440ux43eux433ux440ux430ux43cux43cux43dux43eux433ux43e-ux43aux43eux43cux43fux43bux435ux43aux441ux430}

Проведенный анализ показывает, что существующие на рынке решения делятся
на две категории:

\begin{enumerate}
\def\labelenumi{\arabic{enumi}.}
\item
  \textbf{Инструменты логической виртуализации} (WireMock, MockServer,
  Mountebank) -- отлично справляются с задачей «что ответить на запрос»,
  но не решают задачу управления инфраструктурой (запуск, перезапуск,
  обновление версий приложений).
\item
  \textbf{Серверы приложений} (Apache Tomcat) -- умеют управлять
  жизненным циклом, но слишком тяжеловесны, требуют специфического
  формата упаковки (WAR) и сложны в динамической настройке сетевых
  параметров.
\end{enumerate}

Ни одно из рассмотренных решений не предоставляет функционала
«легковесного агента» для управления сторонними Java-процессами «как
есть» (без перепаковки).

Разработка собственной автоматизированной информационной системы
обоснована следующими факторами:

\begin{enumerate}
\def\labelenumi{\arabic{enumi}.}
\item
  \textbf{Специфика задачи (Lifecycle Management):} Разрабатываемая
  система решает задачу оркестрации \emph{процессов} операционной
  системы. Она позволяет загружать бинарный файл заглушки (JAR),
  полученный от разработчиков, и управлять его исполнением на удаленном
  сервере. Это исключает трудозатраты на пересборку заглушек под
  WireMock или Tomcat.
\item
  \textbf{Эффективность использования ресурсов:} Перенос логики
  управления на \textbf{Python} и использование асинхронного фреймворка
  \textbf{FastAPI} позволяет минимизировать накладные расходы самого
  управляющего контура. Система работает поверх заглушек, не внедряясь в
  их память, в отличие от тяжелых Java-контейнеров.
\item
  \textbf{Централизация и удобство эксплуатации:} Реализация архитектуры
  «Master-Agent» объединяет разрозненные серверы в единый кластер.
  Единая панель управления (Dashboard) позволяет инженеру тестирования
  массово обновлять версии заглушек и менять настройки лимитов (RPS) в
  реальном времени, что невозможно при использовании Standalone-решений.
\item
  \textbf{Технологическая независимость и импортозамещение:} Система
  разрабатывается на базе открытых технологий, адаптирована для работы в
  \textbf{ОС Astra Linux} и не зависит от внешних репозиториев (Maven
  Central, Docker Hub) или лицензионных ключей, что гарантирует
  стабильность работы в закрытых корпоративных контурах РФ.
\end{enumerate}

\subsubsection{Вывод по
подразделу}\label{ux432ux44bux432ux43eux434-ux43fux43e-ux43fux43eux434ux440ux430ux437ux434ux435ux43bux443}

Разработка собственного программного комплекса является единственным
целесообразным решением, закрывающим пробел между инструментами
логического мокирования и средствами управления инфраструктурой,
обеспечивая высокую производительность и гибкость конфигурации тестовых
сред.

\textbf{Список источников (для данного подраздела):}

\begin{enumerate}
\def\labelenumi{\arabic{enumi}.}
\item
  \textbf{WireMock Documentation.} -- Текст: электронный //
  WireMock.org: официальный сайт. -- URL:
  \url{https://wiremock.org/docs/} (дата обращения: 09.01.2025).
\item
  \textbf{MockServer Documentation.} -- Текст: электронный //
  Mock-server.com: официальный сайт. -- URL:
  \url{https://www.mock-server.com/} (дата обращения: 09.01.2025).
\item
  \textbf{Mountebank Documentation.} -- Текст: электронный //
  Mbtest.org: официальный сайт. -- URL:
  \href{https://www.google.com/search?q=http://www.mbtest.org/docs/&authuser=1}{http://www.mbtest.org/docs/}
  (дата обращения: 09.01.2025).
\item
  \textbf{Apache Tomcat 9 Documentation.} -- Текст: электронный //
  Apache.org: официальный сайт. -- URL:
  \url{https://tomcat.apache.org/tomcat-9.0-doc/index.html} (дата
  обращения: 09.01.2025).
\item
  \textbf{Tomcat Clustering/Session Replication How-To.} -- Текст:
  электронный // Apache.org. -- URL:
  \url{https://tomcat.apache.org/tomcat-9.0-doc/cluster-howto.html}
  (дата обращения: 09.01.2025).
\item
  \textbf{Apache Tomcat Configuration Reference: The Rate Limit Filter.}
  -- Текст: электронный // Apache.org. -- URL:
  \href{https://www.google.com/search?q=https://tomcat.apache.org/tomcat-9.0-doc/config/filter.html\%23Rate_Limit_Filter&authuser=1}{https://tomcat.apache.org/tomcat-9.0-doc/config/filter.html\#Rate\_Limit\_Filter}
  (дата обращения: 09.01.2025).
\end{enumerate}

\subsection{Формирование функциональных и технических требований к
системе}\label{ux444ux43eux440ux43cux438ux440ux43eux432ux430ux43dux438ux435-ux444ux443ux43dux43aux446ux438ux43eux43dux430ux43bux44cux43dux44bux445-ux438-ux442ux435ux445ux43dux438ux447ux435ux441ux43aux438ux445-ux442ux440ux435ux431ux43eux432ux430ux43dux438ux439-ux43a-ux441ux438ux441ux442ux435ux43cux435}

На основании проведенного анализа предметной области и специфики задач
нагрузочного тестирования сформированы требования к разрабатываемому
программному комплексу. Система классифицируется как инструмент
автоматизации инфраструктуры нагрузочного тестирования. Архитектура
решения должна быть гибкой, допускающей работу как в распределенном
контуре, так и в рамках одного физического или виртуального сервера.

\subsubsection{Требования к режимам
функционирования}\label{ux442ux440ux435ux431ux43eux432ux430ux43dux438ux44f-ux43a-ux440ux435ux436ux438ux43cux430ux43c-ux444ux443ux43dux43aux446ux438ux43eux43dux438ux440ux43eux432ux430ux43dux438ux44f}

Система должна поддерживать два сценария развертывания без изменения
кодовой базы:

\begin{enumerate}
\def\labelenumi{\arabic{enumi}.}
\item
  \textbf{Распределенный режим (Cluster Mode):} Классическая
  клиент-серверная архитектура, где управляющий узел (Master)
  администрирует множество удаленных узлов (Agents), на которых
  выполняются заглушки. Взаимодействие осуществляется по сети.
\item
  \textbf{Автономный режим (Standalone Mode):} Консолидированный запуск
  всех компонентов системы на одном сервере. В этом режиме управляющий
  модуль и модуль исполнения заглушек функционируют в пределах одной
  операционной системы, обеспечивая изоляцию контура тестирования без
  необходимости сетевого взаимодействия между узлами инфраструктуры.
\end{enumerate}

\subsubsection{Функциональные
требования}\label{ux444ux443ux43dux43aux446ux438ux43eux43dux430ux43bux44cux43dux44bux435-ux442ux440ux435ux431ux43eux432ux430ux43dux438ux44f}

Функциональные требования определяют перечень задач, которые система
должна выполнять для обеспечения процессов эмуляции сервисов.

\paragraph{Требования к подсистеме управления жизненным циклом
(Runtime):}\label{ux442ux440ux435ux431ux43eux432ux430ux43dux438ux44f-ux43a-ux43fux43eux434ux441ux438ux441ux442ux435ux43cux435-ux443ux43fux440ux430ux432ux43bux435ux43dux438ux44f-ux436ux438ux437ux43dux435ux43dux43dux44bux43c-ux446ux438ux43aux43bux43eux43c-runtime}

\begin{itemize}
\item
  \textbf{Поддержка форматов артефактов:} Система должна обеспечивать
  запуск Java-приложений, поставляемых в форматах исполняемых файлов
  .jar и веб-архивов .war.
\item
  \textbf{Управление процессами:} Обеспечение запуска, корректной
  остановки и принудительного завершения процессов заглушек.
\item
  \textbf{Конфигурация запуска:} Возможность задания аргументов
  виртуальной машины (JVM Options), включая параметры памяти (-Xms,
  -Xmx) и системные свойства, через графический интерфейс управляющего
  узла (Master) перед стартом заглушки.
\item
  \textbf{Мониторинг:} Отслеживание статуса процесса (PID) и доступности
  его сетевого порта.
\item
  \textbf{Работа с логами:} Сбор стандартных потоков вывода (stdout,
  stderr) запущенных процессов и их трансляция в интерфейс системы для
  отладки.
\end{itemize}

\paragraph{Требования к подсистеме проксирования запросов (Proxy
Engine):}\label{ux442ux440ux435ux431ux43eux432ux430ux43dux438ux44f-ux43a-ux43fux43eux434ux441ux438ux441ux442ux435ux43cux435-ux43fux440ux43eux43aux441ux438ux440ux43eux432ux430ux43dux438ux44f-ux437ux430ux43fux440ux43eux441ux43eux432-proxy-engine}

\begin{itemize}
\item
  \textbf{Логика обработки запросов:} Система должна работать как
  прозрачный прокси-сервер (Transparent Proxy) на прикладном уровне,
  реализуя следующий алгоритм:

  \begin{enumerate}
  \def\labelenumi{\arabic{enumi}.}
  \item
    Прием входящего HTTP-запроса от внешней системы.
  \item
    Извлечение метаданных: HTTP-метод, заголовки, тело запроса и целевой
    путь (URI).
  \item
    Выполнение нового запроса к целевой заглушке (на её локальный порт)
    с использованием извлеченного метода, заголовков и тела.
  \item
    Получение ответа от заглушки.
  \item
    Возврат полученного ответа инициатору запроса в неизменном виде.
  \end{enumerate}
\item
  \textbf{Управление нагрузкой (Rate Limiting):}

  \begin{itemize}
  \item
    Наличие механизма ограничения количества запросов в секунду (RPS)
    для конкретной заглушки.
  \item
    Активация лимита производится через вызов конфигурационного
    эндпоинта системы (API). Динамическое изменение алгоритма
    ограничения в процессе теста не требуется.
  \item
    При превышении лимита система должна возвращать HTTP-код 429 (Too
    Many Requests), не передавая запрос заглушке.
  \end{itemize}
\end{itemize}

\paragraph{Требования к управляющему модулю
(Master/Dashboard):}\label{ux442ux440ux435ux431ux43eux432ux430ux43dux438ux44f-ux43a-ux443ux43fux440ux430ux432ux43bux44fux44eux449ux435ux43cux443-ux43cux43eux434ux443ux43bux44e-masterdashboard}

\begin{itemize}
\item
  \textbf{Управление артефактами:} Загрузка бинарных файлов (.jar, .war)
  в хранилище системы через веб-интерфейс.
\item
  \textbf{Реестр сервисов:} Отображение списка всех загруженных
  заглушек, их текущего статуса («Запущен», «Остановлен», «Ошибка») и
  адресов.
\item
  \textbf{Интерфейс конфигурации:} Предоставление полей ввода для
  настройки параметров запуска (JVM args) для каждого сервиса.
\end{itemize}

\subsubsection{Нефункциональные (технические)
требования}\label{ux43dux435ux444ux443ux43dux43aux446ux438ux43eux43dux430ux43bux44cux43dux44bux435-ux442ux435ux445ux43dux438ux447ux435ux441ux43aux438ux435-ux442ux440ux435ux431ux43eux432ux430ux43dux438ux44f}

\paragraph{Требования к
производительности}\label{ux442ux440ux435ux431ux43eux432ux430ux43dux438ux44f-ux43a-ux43fux440ux43eux438ux437ux432ux43eux434ux438ux442ux435ux43bux44cux43dux43eux441ux442ux438}

\begin{itemize}
\item
  \textbf{Вносимая задержка (Overhead):} Накладные расходы на
  проксирование (время обработки запроса внутри системы без учета
  времени работы самой заглушки) не должны превышать 20 мс.
\item
  \textbf{Пропускная способность:} Система должна обеспечивать высокую
  производительность проксирующего слоя. Снижение максимальной
  пропускной способности (RPS) заглушки при работе через систему не
  должно превышать 10\% относительно прямого обращения к заглушке.
\end{itemize}

\paragraph{Требования к
надежности}\label{ux442ux440ux435ux431ux43eux432ux430ux43dux438ux44f-ux43a-ux43dux430ux434ux435ux436ux43dux43eux441ux442ux438}

\begin{itemize}
\item
  \textbf{Сохранение конфигурации:} При перезапуске системы (или
  сервера) должны сохраняться настройки запуска заглушек и ссылки на
  загруженные артефакты (Persistence).
\item
  \textbf{Изоляция сбоев:} Ошибка в одной заглушке не должна влиять на
  работоспособность остальных запущенных сервисов и самой управляющей
  системы.
\end{itemize}

\subsubsection{Требования к программно-аппаратному
обеспечению}\label{ux442ux440ux435ux431ux43eux432ux430ux43dux438ux44f-ux43a-ux43fux440ux43eux433ux440ux430ux43cux43cux43dux43e-ux430ux43fux43fux430ux440ux430ux442ux43dux43eux43cux443-ux43eux431ux435ux441ux43fux435ux447ux435ux43dux438ux44e}

\paragraph{Системные
требования}\label{ux441ux438ux441ux442ux435ux43cux43dux44bux435-ux442ux440ux435ux431ux43eux432ux430ux43dux438ux44f}

\begin{itemize}
\item
  \textbf{Операционная система:} Серверная часть должна функционировать
  под управлением \textbf{ОС Astra Linux Special Edition} (версия 1.7 и
  выше) или аналогичных Linux-дистрибутивов.
\item
  \textbf{Среда исполнения Java:} Поддержка \textbf{OpenJDK 17} для
  запуска пользовательских артефактов.
\end{itemize}

\paragraph{Стек
разработки}\label{ux441ux442ux435ux43a-ux440ux430ux437ux440ux430ux431ux43eux442ux43aux438}

\begin{itemize}
\item
  \textbf{Язык программирования:} Python (версия 3.10+) с использованием
  асинхронного подхода.
\item
  \textbf{Веб-фреймворк:} FastAPI (для реализации
  высокопроизводительного ввода-вывода).
\item
  \textbf{СУБД:} Redis (для хранения конфигураций и реализации счетчиков
  Rate Limiter).
\end{itemize}

\subsubsection{Требования к информационной
безопасности}\label{ux442ux440ux435ux431ux43eux432ux430ux43dux438ux44f-ux43a-ux438ux43dux444ux43eux440ux43cux430ux446ux438ux43eux43dux43dux43eux439-ux431ux435ux437ux43eux43fux430ux441ux43dux43eux441ux442ux438}

\begin{itemize}
\item
  \textbf{Контроль доступа:} Доступ к административному интерфейсу и API
  управления должен быть ограничен.
\item
  \textbf{Безопасность данных:} Загружаемые исполняемые файлы должны
  храниться в изолированных директориях с правами доступа, исключающими
  их несанкционированную модификацию третьими лицами.
\end{itemize}

\subsection{}\label{section}

\subsection{Вывод по
разделу}\label{ux432ux44bux432ux43eux434-ux43fux43e-ux440ux430ux437ux434ux435ux43bux443}

В первой главе выпускной квалификационной работы был проведен
комплексный анализ предметной области и обоснована необходимость
создания специализированного инструментария для автоматизации
нагрузочного тестирования.

В ходе исследования выявлено, что использование традиционных серверов
приложений и существующих аналогов (WireMock, MockServer, Mountebank)
для управления распределенной сетью заглушек сопряжено с избыточным
потреблением ресурсов и сложностями в эксплуатации, а также не
удовлетворяет требованиям по импортозамещению.

На основании сравнительного анализа обоснован выбор технологического
стека: язык программирования \textbf{Python} с использованием
асинхронного подхода, веб-фреймворк \textbf{FastAPI} для обеспечения
высокой пропускной способности и In-Memory СУБД \textbf{Redis} для
хранения состояний с минимальной задержкой.

Определена монолитная архитектура будущего решения с управлением
удалёнными хостами через SSH/SCP и сформированы функциональные и
технические требования к системе, включая поддержку форматов .jar и
.war, механизмы ограничения нагрузки (Rate Limiting) и совместимость с
ОС Astra Linux. Полученные результаты служат базисом для проектирования
и технической реализации системы во второй главе.

\section{ПРОЕКТИРОВАНИЕ И ТЕХНИЧЕСКАЯ РЕАЛИЗАЦИЯ СИСТЕМЫ
УПРАВЛЕНИЯ}\label{ux43fux440ux43eux435ux43aux442ux438ux440ux43eux432ux430ux43dux438ux435-ux438-ux442ux435ux445ux43dux438ux447ux435ux441ux43aux430ux44f-ux440ux435ux430ux43bux438ux437ux430ux446ux438ux44f-ux441ux438ux441ux442ux435ux43cux44b-ux443ux43fux440ux430ux432ux43bux435ux43dux438ux44f}

\subsection{Разработка архитектуры распределённого программного
комплекса}\label{ux440ux430ux437ux440ux430ux431ux43eux442ux43aux430-ux430ux440ux445ux438ux442ux435ux43aux442ux443ux440ux44b-ux440ux430ux441ux43fux440ux435ux434ux435ux43bux451ux43dux43dux43eux433ux43e-ux43fux440ux43eux433ux440ux430ux43cux43cux43dux43eux433ux43e-ux43aux43eux43cux43fux43bux435ux43aux441ux430}

\subsubsection{Выбор архитектурного
стиля}\label{ux432ux44bux431ux43eux440-ux430ux440ux445ux438ux442ux435ux43aux442ux443ux440ux43dux43eux433ux43e-ux441ux442ux438ux43bux44f}

Проектирование архитектуры является фундаментальным этапом разработки
программного комплекса: именно на данной стадии закладываются
структурные свойства системы, определяющие её масштабируемость,
надёжность, сопровождаемость и производительность в долгосрочной
перспективе. В рамках настоящего раздела проводится анализ применимых
архитектурных подходов, обосновывается выбор итогового архитектурного
решения и формализуется структура разработанного программного комплекса.

В современной инженерии программного обеспечения наибольшее
распространение получили два конкурирующих архитектурных стиля:
монолитная и микросервисная архитектуры. Для обоснованного выбора
необходимо рассмотреть их характеристики в контексте требований,
сформулированных в первой главе настоящей работы.

\textbf{Микросервисная архитектура} предполагает декомпозицию системы на
множество независимо разворачиваемых сервисов, каждый из которых
реализует конкретную бизнес-функцию и взаимодействует с остальными
посредством сетевых протоколов -- как правило, HTTP REST или gRPC.
Данный подход демонстрирует преимущества в условиях работы крупных
распределённых команд разработчиков, при необходимости независимого
масштабирования отдельных компонентов и при высокой частоте выпуска
обновлений изолированных частей системы. К характерным достоинствам
микросервисной архитектуры относятся: независимость технологического
стека отдельных сервисов, устойчивость к отказам на уровне изолированных
процессов, а также возможность горизонтального масштабирования
конкретных узлов без перезапуска всей системы.

Вместе с тем микросервисный подход сопряжён с существенными
эксплуатационными издержками, принципиально важными в контексте
настоящего проекта. Во-первых, каждый сервис требует собственной
инфраструктуры развёртывания: отдельный сетевой порт, конфигурационный
файл, процесс инициализации и мониторинга. Во-вторых, взаимодействие
сервисов по сети порождает дополнительные задержки и накладные расходы
на сериализацию и десериализацию данных при каждом межсервисном вызове.
В-третьих, полноценная реализация микросервисной архитектуры требует
развёртывания дополнительных инфраструктурных компонентов: системы
оркестрации контейнеров, сервисного реестра и системы трассировки
запросов. Перечисленные компоненты существенно увеличивают сложность
системы и создают зависимости от зарубежного программного обеспечения,
что противоречит требованию технологической независимости,
сформулированному в первой главе.

\textbf{Монолитная архитектура} объединяет все компоненты системы в
рамках единого процесса операционной системы, разделяя общее адресное
пространство памяти, единую точку входа и общую конфигурацию
развёртывания. Несмотря на устойчивое представление о монолите как об
устаревшем архитектурном решении, данный подход демонстрирует высокую
эффективность в условиях, характерных для разрабатываемого комплекса:
команда разработки состоит из одного специалиста, функциональные
требования чётко определены и стабильны, а потребности в раздельном
масштабировании отдельных внутренних компонентов отсутствуют.

Ключевым аргументом в пользу монолитного подхода является специфика
решаемой задачи: разрабатываемая система является «тонким»
координирующим слоем над управляемыми имитационными сервисами. Её
основная функция -- оркестрация процессов операционной системы на
целевых хостах, маршрутизация входящего трафика и хранение
конфигурационных данных. Введение внутрисистемных сетевых вызовов в
рамках гипотетической микросервисной декомпозиции создало бы накладные
расходы, несоразмерные выигрышу от разделения компонентов.

Распределённость системы, заявленная в названии раздела, реализуется не
на уровне внутреннего устройства приложения, а на уровне топологии
развёртывания: единый управляющий процесс координирует целевые хосты
посредством защищённых сетевых протоколов. Таким образом,
распределённость является внешней характеристикой инфраструктуры, а не
принципом внутренней декомпозиции.

На основании проведённого анализа для реализации программного комплекса
выбрана монолитная архитектура с явно выраженной внутренней модульной
структурой. Данное решение обеспечивает минимальные накладные расходы на
взаимодействие компонентов, простоту развёртывания и эксплуатации в
закрытом корпоративном контуре, технологическую независимость от внешних
оркестраторов, а также высокую сопровождаемость и читаемость единой
кодовой базы.

В таблице 2 приведён сравнительный анализ двух рассмотренных
архитектурных подходов по критериям, наиболее значимым для данного
проекта.

\emph{Таблица 2 -- Сравнение архитектурных подходов}

{\def\LTcaptype{none} % do not increment counter
\begin{longtable}[]{@{}
  >{\raggedright\arraybackslash}p{(\linewidth - 4\tabcolsep) * \real{0.2979}}
  >{\raggedright\arraybackslash}p{(\linewidth - 4\tabcolsep) * \real{0.3511}}
  >{\raggedright\arraybackslash}p{(\linewidth - 4\tabcolsep) * \real{0.3511}}@{}}
\toprule\noalign{}
\begin{minipage}[b]{\linewidth}\centering
\textbf{Критерий}
\end{minipage} & \begin{minipage}[b]{\linewidth}\centering
\textbf{Микросервисная архитектура}
\end{minipage} & \begin{minipage}[b]{\linewidth}\centering
\textbf{Монолитная архитектура (выбрано)}
\end{minipage} \\
\midrule\noalign{}
\endhead
\bottomrule\noalign{}
\endlastfoot
Сложность развёртывания & Высокая (оркестратор, реестр, трассировка) &
Низкая (один процесс, один конфиг) \\
Накладные расходы на вызовы & Сетевые задержки между сервисами &
Отсутствуют (in-process вызовы) \\
Изоляция отказов & На уровне процессов ОС & На уровне модулей
(программная) \\
Потребность в доп. инфраструктуре & Docker/K8s, Consul, Jaeger & Только
Python + Redis \\
Подходит команде из 1 человека & Нет (высокий DevOps-оверхед) & Да
(минимальные операционные издержки) \\
Соответствие требованиям импортозамещения & Частичное (зависимость от
K8s) & Полное (открытые технологии) \\
\end{longtable}
}

\subsubsection{Топология развёртывания и гибкость
инфраструктуры}\label{ux442ux43eux43fux43eux43bux43eux433ux438ux44f-ux440ux430ux437ux432ux451ux440ux442ux44bux432ux430ux43dux438ux44f-ux438-ux433ux438ux431ux43aux43eux441ux442ux44c-ux438ux43dux444ux440ux430ux441ux442ux440ux443ux43aux442ux443ux440ux44b}

Одним из ключевых проектных решений, определивших структуру системы,
является принцип гибкого размещения имитационных сервисов. В отличие от
систем, жёстко делящихся на «локальный» и «кластерный» режимы,
разработанный программный комплекс предоставляет оператору возможность
назначать произвольный целевой хост для каждой заглушки независимо.

Распределённость системы реализуется посредством реестра хостов: каждый
зарегистрированный хост описывается своим сетевым адресом, учётной
записью для подключения и рабочей директорией. При развёртывании
конкретной заглушки оператор выбирает один из хостов реестра в качестве
целевого. Это позволяет в рамках единого экземпляра управляющего
приложения одновременно:

\begin{itemize}
\item
  размещать часть заглушек непосредственно на сервере, где работает
  управляющее приложение (адрес 127.0.0.1 или hostname сервера), -- что
  соответствует традиционному «локальному» сценарию;
\item
  размещать другие заглушки на удалённых серверах через SSH/SCP -- что
  соответствует «кластерному» сценарию;
\item
  совмещать оба варианта в рамках одной инфраструктуры без каких-либо
  изменений в конфигурации приложения.
\end{itemize}

Таким образом, граница между «локальным» и «распределённым» режимами
определяется не конфигурацией запуска приложения, а составом реестра
хостов и выбором оператора при регистрации каждой заглушки. Такой подход
значительно увеличивает гибкость эксплуатации системы и позволяет плавно
наращивать инфраструктуру без вмешательства в работу уже запущенных
сервисов.

С точки зрения внутренней архитектуры это означает, что модуль
управления жизненным циклом всегда работает по единому алгоритму,
независимо от того, является ли целевой хост локальным или удалённым.
При работе с localhost SSH-соединение устанавливается к локальному
SSH-серверу по стандартному порту 22, что унифицирует код и устраняет
необходимость в двух ветках логики для локального и удалённого
сценариев.

\subsubsection{Контекстная диаграмма
системы}\label{ux43aux43eux43dux442ux435ux43aux441ux442ux43dux430ux44f-ux434ux438ux430ux433ux440ux430ux43cux43cux430-ux441ux438ux441ux442ux435ux43cux44b}

Для формализованного описания архитектуры программного комплекса
применяется нотация C4 Model, предусматривающая четыре уровня
детализации: контекстный (Level 1), контейнерный (Level 2), компонентный
(Level 3) и кодовый (Level 4). Нотация обеспечивает однозначность
интерпретации архитектурных решений и соответствует требованиям
методических указаний по применению стандартизованных схем для
представления структуры системы.

Контекстная диаграмма (рисунок 1) отображает систему как единый «чёрный
ящик» в окружении внешних акторов и смежных систем. На данном уровне
абстракции детали внутреннего устройства системы не раскрываются;
основное внимание уделяется определению границ системы и характеру её
взаимодействия с внешним окружением.

\includesvg[width=6.49653in,height=4.27778in]{media/image2.svg}\\
Рисунок 1 -- контекстная диаграмма системы

Внешними акторами системы являются:

\begin{itemize}
\item
  инженер тестирования -- специалист, осуществляющий управление
  инфраструктурой имитационных сервисов через веб-интерфейс системы;
  выполняет операции регистрации артефактов, запуска и остановки
  заглушек, настройки параметров и просмотра метрик;
\item
  потребитель -- любой HTTP-клиент, направляющий запросы к имитационным
  сервисам через прокси-движок системы: это может быть инструмент
  нагрузочного тестирования (Apache JMeter, Gatling и др.), тестируемая
  система (System Under Test) или иной HTTP-клиент;
\item
  Prometheus -- внешняя система сбора и хранения метрик, периодически
  опрашивающая специализированный эндпоинт системы для получения
  телеметрических данных;
\item
  целевой хост -- сервер под управлением ОС Astra Linux, на котором
  непосредственно выполняются Java-процессы имитационных сервисов; может
  совпадать с сервером управляющего приложения.
\end{itemize}

workspace \{

model \{

engineer = person "Инженер тестирования" \{

description "Управляет инфраструктурой заглушек через веб-интерфейс"

\}

consumer = person "Потребитель" \{

description "Любой клиент, направляющий HTTP-запросы к заглушкам:
нагрузочный инструмент, тестируемая система или иной HTTP-клиент"

\}

prometheus = softwareSystem "Prometheus" \{

description "Внешняя система сбора и хранения метрик"

\}

mockSystem = softwareSystem "Программный комплекс управления заглушками"
\{

description "Монолитное Python/FastAPI-приложение"

webUI = container "Веб-интерфейс (Jinja2 SSR)" \{

description "HTML-страницы, формируемые на стороне сервера: Dashboard,
Настройки, Логи"

technology "FastAPI / Jinja2 / HTML"

\}

apiLayer = container "REST API" \{

description "Бизнес-логика: управление заглушками, проксирование,
настройки"

technology "Python 3.10+ / FastAPI / Uvicorn (ASGI)"

lifecycleMgr = component "Lifecycle Manager" \{

description "Регистрация (загрузка + SCP на хост), запуск, остановка,
удаление Java-процессов"

technology "Python asyncio / paramiko (SSH/SCP)"

\}

proxyEngine = component "Proxy Engine" \{

description "Прозрачная маршрутизация HTTP-запросов от потребителей к
целевым заглушкам"

technology "Python httpx (async)"

\}

rateLimiter = component "Rate Limiter" \{

description "Алгоритм Fixed Window Counter: ограничение RPS на уровне
каждой заглушки"

technology "Redis INCR + EXPIRE"

\}

settingsModule = component "Settings Module" \{

description "Управление учётными записями SSH и реестром хостов
развёртывания"

technology "Python / Redis"

\}

cryptoService = component "Crypto Service" \{

description "Шифрование/дешифрование паролей SSH перед сохранением в
Redis"

technology "Python cryptography (Fernet)"

\}

metricsExporter = component "Metrics Exporter" \{

description "Эндпоинт /metrics в формате Prometheus exposition format"

technology "Python"

\}

\}

redis = container "Redis" \{

description "In-Memory хранилище конфигураций и состояний"

technology "Redis 7.x (AOF persistence)"

\}

\}

targetHost = softwareSystem "Целевой хост" \{

description "Сервер ОС Astra Linux с запущенными заглушками. Может
совпадать с сервером управляющего приложения."

\}

engineer -\textgreater{} webUI "Управляет через браузер (HTTP)"

consumer -\textgreater{} proxyEngine "HTTP-запросы к заглушкам"

prometheus -\textgreater{} apiLayer "Scrape /metrics (HTTP GET)"

webUI -\textgreater{} apiLayer "Вызовы REST API (HTTP)"

apiLayer -\textgreater{} redis "Чтение/запись конфигураций и состояний"

lifecycleMgr -\textgreater{} targetHost "SCP артефакта при регистрации;
SSH: запуск/остановка/удаление процесса и файла"

proxyEngine -\textgreater{} targetHost "Перенаправление HTTP-запросов к
работающей заглушке"

settingsModule -\textgreater{} cryptoService "Шифрование пароля перед
сохранением; дешифрование перед использованием"

rateLimiter -\textgreater{} redis "Инкремент счётчика и проверка лимита"

\}

views \{

systemContext mockSystem "SystemContext" \{

include *

autolayout lr

title "Рисунок 2.1 -- Контекстная диаграмма системы (C4 Level 1)"

\}

container mockSystem "Containers" \{

include *

autolayout lr

title "Рисунок 2.2 -- Контейнерная диаграмма (C4 Level 2)"

\}

component apiLayer "Components" \{

include *

autolayout tb

title "Рисунок 2.3 -- Компонентная диаграмма REST API (C4 Level 3)"

\}

styles \{

element "Person" \{

shape Person

background \#08427B

color \#ffffff

\}

element "Software System" \{

background \#1168BD

color \#ffffff

\}

element "Container" \{

background \#438DD5

color \#ffffff

\}

element "Component" \{

background \#85BBF0

color \#000000

\}

\}

\}

\}

\subsubsection{Контейнерная
диаграмма}\label{ux43aux43eux43dux442ux435ux439ux43dux435ux440ux43dux430ux44f-ux434ux438ux430ux433ux440ux430ux43cux43cux430}

Контейнерная диаграмма (рисунок 2) раскрывает внутреннее устройство
программного комплекса на уровне самостоятельно исполняемых программных
единиц -- контейнеров. В терминологии C4 Model «контейнер» обозначает
логически обособленную программную единицу, а не технологию
Docker-контейнеризации. Для каждого контейнера указывается его
технологический стек и зона ответственности.

\includesvg[width=6.49653in,height=4.15625in]{media/image4.svg}\\
Рисунок 2 -- контейнерная диаграмма системы

В составе программного комплекса выделяются следующие контейнеры.

\textbf{Веб-интерфейс (Jinja2 SSR).} Реализует пользовательский
интерфейс на основе технологии Server-Side Rendering с использованием
шаблонизатора Jinja2, встроенного в экосистему FastAPI. HTML-страницы
формируются на стороне сервера и передаются браузеру в готовом виде.
Данный подход минимизирует объём JavaScript-кода на стороне клиента и
обеспечивает полную функциональность интерфейса без зависимости от
внешних CDN. Веб-интерфейс включает: главную панель (Dashboard) со
списком заглушек, страницу управления хостами и настроек системы,
страницу просмотра журналов процессов.

\textbf{REST API (FastAPI).} Центральный контейнер системы, реализующий
всю бизнес-логику. Запускается под управлением ASGI-сервера Uvicorn.
REST API предоставляет конечные точки для взаимодействия с
веб-интерфейсом, а также программный интерфейс для возможной интеграции
со сторонними системами автоматизации. Внутренняя структура контейнера
детализируется на компонентном уровне в разделе 2.1.5.

\textbf{Redis.} In-Memory хранилище структур данных, выполняющее функции
постоянного хранилища конфигурационных данных системы и оперативного
хранилища счётчиков алгоритма ограничения запросов. Redis функционирует
в режиме персистентности AOF (Append Only File), что обеспечивает
сохранность данных при аварийном перезапуске управляющей системы.
Детальное описание структур данных Redis приводится в разделе 2.2.

Система намеренно не содержит выделенного контейнера хранилища
артефактов. Архитектурное решение «1 заглушка -- 1 хост» исключает
необходимость повторного развёртывания одного артефакта на несколько
хостов, а значит, постоянное хранение загруженного файла на сервере
управляющего приложения избыточно. При регистрации заглушки её бинарный
файл принимается из браузера во временный каталог, немедленно копируется
на целевой хост посредством SCP и удаляется с сервера управляющего
приложения. Единственной постоянной копией артефакта остаётся файл на
целевом хосте в рабочей директории, указанной в реестре хостов.

\subsubsection{Компонентная диаграмма REST
API}\label{ux43aux43eux43cux43fux43eux43dux435ux43dux442ux43dux430ux44f-ux434ux438ux430ux433ux440ux430ux43cux43cux430-rest-api}

Компонентная диаграмма (рисунок 3) детализирует внутреннее устройство
контейнера REST API, раскрывая составляющие его программные модули --
компоненты.

\includesvg[width=6.49653in,height=3.9875in]{media/image6.svg}\\
Рисунок 3 -- компонентная диаграмма системы

\textbf{Lifecycle Manager.} Реализует операции управления жизненным
циклом имитационных сервисов. При регистрации новой заглушки модуль
принимает загруженный бинарный файл во временный каталог на сервере
управляющего приложения, копирует его на целевой хост посредством SCP в
рабочую директорию, указанную в записи хоста, и немедленно удаляет
временный файл. Путь к файлу на целевом хосте не сохраняется в Redis --
он вычисляется при необходимости как конкатенация поля working\_dir из
записи hosts:\{hostname\} и имени файла заглушки. При запуске модуль
устанавливает SSH-соединение с целевым хостом, определяет свободный порт
в заданном диапазоне и инициирует Java-процесс с указанными
JVM-аргументами в фоновом режиме (nohup). При остановке выполняется
корректное завершение процесса через SSH. При удалении заглушки
производятся: остановка процесса (если запущен), удаление бинарного
файла с целевого хоста командой SSH rm и удаление записи из Redis.

\textbf{Proxy Engine.} Реализует функциональность прозрачного обратного
прокси-сервера на прикладном уровне. При поступлении HTTP-запроса движок
извлекает адрес и порт целевой заглушки из Redis, формирует новый
HTTP-запрос с сохранением метода, заголовков и тела оригинального
запроса, а затем возвращает ответ заглушки инициатору в неизменном виде.
Для асинхронного выполнения исходящих запросов используется библиотека
httpx.

\textbf{Rate Limiter.} Реализует алгоритм «Счётчик с фиксированным
временным окном» (Fixed Window Counter). Компонент функционирует на
уровне каждой заглушки независимо и управляется флагом
rate\_limit\_enabled, переключаемым оператором через Dashboard без
остановки заглушки. Детальное описание алгоритма приводится в разделе
2.4.

\textbf{Settings Module.} Обеспечивает управление конфигурационными
параметрами системы -- учётными записями ОС Linux и реестром хостов
развёртывания. Подробное описание приведено в разделе 2.1.6.

\textbf{Crypto Service.} Выполняет шифрование паролей SSH перед записью
в Redis и дешифрование перед использованием для установки соединения.
Применяется алгоритм AES-128 (Fernet). Ключ шифрования хранится на
файловой системе сервера в файле с правами 600.

\textbf{Metrics Exporter.} Реализует эндпоинт /metrics в формате
Prometheus exposition format. Агрегирует данные о RPS, срабатываниях
Rate Limiter, статусах процессов и доступности хостов. Детальное
описание приведено в разделе 2.6.

\subsubsection{Архитектура модуля настроек
системы}\label{ux430ux440ux445ux438ux442ux435ux43aux442ux443ux440ux430-ux43cux43eux434ux443ux43bux44f-ux43dux430ux441ux442ux440ux43eux435ux43a-ux441ux438ux441ux442ux435ux43cux44b}

Модуль настроек (Settings Module) предоставляет администратору системы
веб-интерфейс для конфигурирования двух категорий параметров: учётных
записей операционной системы Linux, применяемых для SSH/SCP-операций, и
реестра хостов, доступных для развёртывания имитационных сервисов.

\textbf{Управление учётными записями ОС Linux.} Данная категория
позволяет зарегистрировать именованные учётные записи пользователей
операционной системы Linux. Каждая учётная запись используется для
аутентификации при установке SSH-соединения с целевым хостом и включает
имя пользователя (login) и пароль. Применение отдельных учётных записей
с ограниченными правами доступа обусловлено требованиями информационной
безопасности: рабочая директория заглушек должна быть доступна только
указанному пользователю, что минимизирует риски при нештатном поведении
Java-процессов.

\textbf{Механизм защищённого хранения паролей.} Поскольку пароли учётных
записей SSH хранятся в Redis, необходимо обеспечить их
конфиденциальность в случае несанкционированного доступа к базе данных.
Для этого применяется симметричное шифрование на основе алгоритма
AES-128 в режиме CBC, реализованного в библиотеке cryptography пакетом
Fernet. Схема работы следующая:

\begin{itemize}
\item
  при первом запуске приложения генерируется или загружается из
  защищённого конфигурационного файла ключ шифрования (Fernet key),
  который хранится вне Redis -- в файле с правами доступа 600, доступном
  только пользователю, от имени которого запущено приложение;
\item
  перед сохранением пароля в Redis компонент Crypto Service шифрует его
  с использованием Fernet-ключа и сохраняет в виде строки
  Base64-кодированного зашифрованного блока;
\item
  перед использованием пароля для установки SSH-соединения Crypto
  Service выполняет дешифрование, получая исходный пароль в открытом
  виде только в оперативной памяти процесса на время установки
  соединения.
\end{itemize}

Такой подход гарантирует, что даже при получении злоумышленником дампа
Redis пароли останутся недоступными без Fernet-ключа, хранящегося на
файловой системе сервера.

\textbf{Реестр хостов развёртывания.} Данная категория настроек
представляет собой список серверов, доступных для развёртывания
имитационных сервисов. Каждая запись реестра описывает конкретный
целевой хост и содержит параметры, перечисленные в таблице 2.2.

В таблице 3 приведено описание параметров модуля настроек с указанием
типов данных и ограничений.

\emph{Таблица 3 -- Параметры модуля настроек системы}

{\def\LTcaptype{none} % do not increment counter
\begin{longtable}[]{@{}
  >{\raggedright\arraybackslash}p{(\linewidth - 8\tabcolsep) * \real{0.1940}}
  >{\raggedright\arraybackslash}p{(\linewidth - 8\tabcolsep) * \real{0.1381}}
  >{\raggedright\arraybackslash}p{(\linewidth - 8\tabcolsep) * \real{0.1145}}
  >{\raggedright\arraybackslash}p{(\linewidth - 8\tabcolsep) * \real{0.1958}}
  >{\raggedright\arraybackslash}p{(\linewidth - 8\tabcolsep) * \real{0.3576}}@{}}
\toprule\noalign{}
\begin{minipage}[b]{\linewidth}\centering
\textbf{Параметр}
\end{minipage} & \begin{minipage}[b]{\linewidth}\centering
\textbf{Раздел}
\end{minipage} & \begin{minipage}[b]{\linewidth}\centering
\textbf{Тип}
\end{minipage} & \begin{minipage}[b]{\linewidth}\centering
\textbf{Обязательный}
\end{minipage} & \begin{minipage}[b]{\linewidth}\centering
\textbf{Описание}
\end{minipage} \\
\midrule\noalign{}
\endhead
\bottomrule\noalign{}
\endlastfoot
username & Учётная запись & String & Да & Имя пользователя ОС Linux для
SSH-аутентификации \\
password & Учётная запись & String & Да & Пароль (хранится в
зашифрованном виде через Fernet) \\
ssh\_port & Хост & Integer & Нет & Порт SSH (по умолчанию: 22) \\
account\_id & Хост & UUID & Да & Идентификатор привязанной учётной
записи \\
working\_dir & Хост & String (path) & Да & Рабочая директория на хосте
для хранения файлов заглушек \\
java\_path & Хост & String (path) & Да & Абсолютный путь до исполняемого
файла java \\
mock\_port\_min & Хост & Integer & Нет & Нижняя граница диапазона портов
для заглушек (по умолч.: 8100) \\
mock\_port\_max & Хост & Integer & Нет & Верхняя граница диапазона
портов для заглушек (по умолч.: 8200) \\
description & Хост & String & Нет & Произвольный комментарий к хосту
(назначение, окружение) \\
\end{longtable}
}

Параметр java\_path введён в реестр хостов по причине неоднородности
целевой инфраструктуры: в корпоративных средах переменная окружения
JAVA\_HOME нередко не определена или указывает на версию JRE, не
совместимую с требованиями конкретной заглушки. Указание явного пути к
исполняемому файлу Java устраняет неоднозначность и гарантирует запуск
процесса с нужной версией JDK вне зависимости от конфигурации окружения
на хосте.

Параметры mock\_port\_min и mock\_port\_max задают диапазон TCP-портов,
в пределах которого система будет динамически выбирать свободный порт
при запуске каждой новой заглушки на данном хосте. Ограничение диапазона
необходимо для совместимости с политиками межсетевого экрана, которые в
закрытых корпоративных сетях, как правило, разрешают входящий трафик
только на заранее согласованные порты.

Поле description позволяет операторам документировать назначение каждого
хоста непосредственно в интерфейсе системы, что упрощает эксплуатацию
при работе с несколькими окружениями.

Хранение информации о заглушках в Redis организовано в формате
JSON-строки, где ключом служит имя файла заглушки (без пути, например
payment-service-stub.jar). Имя файла является уникальным идентификатором
заглушки в системе: при попытке зарегистрировать файл с именем, уже
присутствующим в Redis, система возвращает ошибку и предлагает оператору
переименовать загружаемый артефакт или удалить существующую запись.
Такой подход упрощает поиск конфигурации по известному имени файла --
операция сводится к одному прямому обращению GET к Redis без
дополнительного индекса. Структура значения имеет следующий вид:

\emph{Пример структуры записи о заглушке в Redis}

"payment-service-stub.jar": \{

"hostname": "test-server-01",

"port": 8105,

"pid": 24731,

"status": "RUNNING",

"jvm\_args": "-Xms128m -Xmx256m",

"rate\_limit": 500,

"rate\_limit\_enabled": false

\}

Путь к бинарному файлу в записи не хранится: он вычисляется при
необходимости как \{working\_dir\}/\{filename\}, где working\_dir
берётся из записи хоста, а filename -- из ключа Redis.

Поле rate\_limit\_enabled является булевым флагом, управляющим
активностью механизма ограничения запросов для конкретной заглушки.
Значение по умолчанию при регистрации новой заглушки -- false (Rate
Limiter отключён). Оператор может изменить значение флага в любой момент
через переключатель на Dashboard без остановки заглушки: изменение
немедленно записывается в Redis и применяется Proxy Engine при обработке
следующего входящего запроса.

Аналогичная JSON-структура применяется для хранения записей о хостах и
учётных записях. Подробное описание всех ключей и структур данных Redis
приводится в разделе 2.2 настоящей главы.

\subsubsection{Диаграмма
развёртывания}\label{ux434ux438ux430ux433ux440ux430ux43cux43cux430-ux440ux430ux437ux432ux451ux440ux442ux44bux432ux430ux43dux438ux44f}

Диаграмма развёртывания (рисунок 4) описывает физическое размещение
компонентов системы на серверной инфраструктуре. Сервер управляющего
приложения содержит два постоянных компонента: FastAPI-процесс под
управлением Uvicorn и экземпляр Redis. Бинарные файлы заглушек постоянно
хранятся исключительно на целевых хостах в рабочих директориях.
Временный каталог на сервере управляющего приложения используется лишь
на время передачи файла (в течение одной SCP-операции) и очищается
немедленно после её завершения.

\includesvg[width=6.49653in,height=1.4125in]{media/image8.svg}\\
Рисунок 4 -- Диаграмма развёртывания

@startuml deployment\_diagram

skinparam backgroundColor \#FFFFFF

skinparam defaultFontName Arial

skinparam defaultFontSize 12

skinparam node \{

BackgroundColor \#D6E4F0

BorderColor \#1A6FA8

BorderThickness 2

\}

skinparam rectangle \{

BackgroundColor \#FAFAFA

BorderColor \#5D6D7E

\}

node "**Сервер управляющего приложения**\textbackslash n//Astra Linux SE
1.7 · Python 3.10 · Uvicorn · Redis 7.x//" as Master \{

rectangle "**«FastAPI-процесс (Uvicorn)»**\textbackslash nREST API +
Веб-интерфейс\textbackslash n+ все модули системы" as FastAPI

rectangle "**«Redis»**\textbackslash nIn-Memory
СУБД\textbackslash nAOF-персистентность" as Redis

FastAPI -\/-\textgreater{} Redis :
Чтение/запись\textbackslash nсостояний (TCP 6379)

\}

node "**«localhost»**\textbackslash n//Тот же физический сервер ·
OpenJDK 17//" as LocalHost \#D5F5E3 \{

rectangle "**«Java-процессы заглушек»**\textbackslash nОдин или
несколько\textbackslash n.jar/.war процессов" as LocalMocks \#FAFAFA

\}

node "**«Удалённый хост 1»**\textbackslash n//Astra Linux SE 1.7 · SSH ·
OpenJDK 17//" as Remote1 \#FEF9E7 \{

rectangle "**«Java-процессы заглушек»**\textbackslash nОдин или
несколько\textbackslash n.jar/.war процессов" as RemoteMocks1 \#FAFAFA

\}

node "**«Удалённый хост 2»**\textbackslash n//Astra Linux SE 1.7 · SSH ·
OpenJDK 17//" as Remote2 \#FEF9E7 \{

rectangle "**«Java-процессы заглушек»**\textbackslash nОдин или
несколько\textbackslash n.jar/.war процессов" as RemoteMocks2 \#FAFAFA

\}

FastAPI -\/-\textgreater{} LocalMocks : Управление процессами
(localhost)

FastAPI -\/-\textgreater{} RemoteMocks1 : SSH/SCP при
регистрации;\textbackslash nHTTP-проксирование запросов

FastAPI -\/-\textgreater{} RemoteMocks2 : SSH/SCP при
регистрации;\textbackslash nHTTP-проксирование запросов

@enduml

На целевых хостах не требуется установки каких-либо компонентов
разрабатываемой системы. Обязательные условия для включения хоста в
реестр: работающий SSH-сервер, настроенная учётная запись пользователя с
правами на запись в рабочую директорию и наличие исполняемого файла Java
по пути, указанному в параметре java\_path. Бинарные файлы заглушек
хранятся в рабочей директории хоста на протяжении всего времени
существования записи о заглушке в системе и удаляются только при явном
удалении заглушки оператором.

\subsubsection{Сценарий взаимодействия
компонентов}\label{ux441ux446ux435ux43dux430ux440ux438ux439-ux432ux437ux430ux438ux43cux43eux434ux435ux439ux441ux442ux432ux438ux44f-ux43aux43eux43cux43fux43eux43dux435ux43dux442ux43eux432}

Для иллюстрации сквозного взаимодействия компонентов рассмотрим пять
типовых сценариев. Принципиальное архитектурное решение: копирование
бинарного артефакта на целевой хост выполняется однократно при
регистрации заглушки; все последующие операции (запуск, остановка,
переключение Rate Limiter) работают с файлом, уже находящимся на целевом
хосте.

\textbf{Сценарий 1. Регистрация новой заглушки.}

Шаг 1. Инженер тестирования загружает .jar- или .war-файл через
веб-интерфейс и задаёт параметры: JVM-аргументы, целевой хост из
реестра, пороговое значение Rate Limiter. Флаг rate\_limit\_enabled
устанавливается в false по умолчанию.

Шаг 2. REST API сохраняет файл во временный каталог на сервере
управляющего приложения.

Шаг 3. Lifecycle Manager получает параметры целевого хоста из Redis
(адрес, учётная запись, рабочая директория). Crypto Service дешифрует
пароль учётной записи.

Шаг 4. Lifecycle Manager копирует файл с сервера управляющего приложения
на целевой хост посредством SCP в рабочую директорию.

Шаг 5. После успешного завершения SCP временный файл удаляется с сервера
управляющего приложения.

Шаг 6. В Redis создаётся JSON-запись заглушки с ключом
mocks:\{filename\}. Имя файла добавляется в реестр mocks:registry.
Статус устанавливается в REGISTERED.

Шаг 7. Веб-интерфейс отображает новую заглушку на Dashboard со статусом
«Зарегистрирована» и переключателем Rate Limiter в положении «Выкл.».

\textbf{Сценарий 2. Запуск зарегистрированной заглушки.}

Шаг 1. Инженер тестирования нажимает кнопку «Запустить» для выбранной
заглушки на Dashboard.

Шаг 2. Браузер отправляет POST-запрос к эндпоинту
/api/v1/mocks/\{mock\_name\}/start.

Шаг 3. Lifecycle Manager считывает из Redis конфигурацию заглушки
(JVM-аргументы) и параметры хоста (адрес, java\_path, учётная запись).
Crypto Service дешифрует пароль.

Шаг 4. Lifecycle Manager определяет свободный TCP-порт в диапазоне
mock\_port\_min / mock\_port\_max через SSH-команду на целевом хосте.

Шаг 5. Lifecycle Manager выполняет SSH-команду: \{java\_path\}
\{jvm\_args\} -jar \{artifact\_path\} -\/-server.port=\{port\}. Процесс
запускается в фоновом режиме (nohup).

Шаг 6. PID процесса, назначенный порт и статус RUNNING сохраняются в
запись заглушки в Redis.

Шаг 7. Proxy Engine начинает маршрутизировать входящие запросы на
зафиксированный порт заглушки.

\textbf{Сценарий 3. Обработка входящего запроса через прокси.}

Шаг 1. Потребитель направляет HTTP-запрос на адрес вида
/\{mock\_name\}/\{path\}.

Шаг 2. Proxy Engine извлекает из Redis запись заглушки по mock\_name:
адрес, порт, статус, флаг rate\_limit\_enabled.

Шаг 3. Если статус отличен от RUNNING -- возвращается HTTP 503 (Service
Unavailable).

Шаг 4а. Если rate\_limit\_enabled = false -- запрос немедленно
передаётся целевой заглушке, счётчики не обновляются.

Шаг 4б. Если rate\_limit\_enabled = true -- Rate Limiter инкрементирует
счётчик rate:\{mock\_name\}:\{window\} и сравнивает с порогом
rate\_limit. При превышении -- возвращается HTTP 429 (Too Many
Requests).

Шаг 5. Proxy Engine асинхронно передаёт запрос целевой заглушке через
httpx и возвращает ответ потребителю в неизменном виде.

\textbf{Сценарий 4. Переключение Rate Limiter на Dashboard.}

Шаг 1. Оператор переключает тумблер Rate Limiter для выбранной заглушки
на Dashboard.

Шаг 2. Браузер отправляет PATCH-запрос к
/api/v1/mocks/\{mock\_name\}/rate-limit с телом \{ "enabled": true/false
\}.

Шаг 3. Settings Module обновляет поле rate\_limit\_enabled в JSON-записи
заглушки в Redis последовательностью операций GET → modify → SET. Для
предотвращения состояния гонки при конкурентных изменениях операция
выполняется с использованием механизма оптимистичной блокировки Redis
(WATCH/MULTI/EXEC).

Шаг 4. Начиная со следующего запроса к данной заглушке, Proxy Engine
применяет обновлённое значение флага без перезапуска.

\textbf{Сценарий 5. Удаление заглушки.}

Шаг 1. Инженер тестирования инициирует удаление заглушки через
веб-интерфейс.

Шаг 2. Если заглушка в состоянии RUNNING -- Lifecycle Manager выполняет
остановку процесса через SSH (SIGTERM, при необходимости SIGKILL).

Шаг 3. Lifecycle Manager удаляет бинарный файл заглушки с целевого хоста
командой SSH rm \{artifact\_path\}.

Шаг 4. Запись о заглушке и все связанные ключи (счётчики Rate Limiter)
удаляются из Redis командами DEL и SREM.

Шаг 5. Запись исчезает из Dashboard.

\subsection{Проектирование структур данных для хранения конфигураций и
состояний в
Redis}\label{ux43fux440ux43eux435ux43aux442ux438ux440ux43eux432ux430ux43dux438ux435-ux441ux442ux440ux443ux43aux442ux443ux440-ux434ux430ux43dux43dux44bux445-ux434ux43bux44f-ux445ux440ux430ux43dux435ux43dux438ux44f-ux43aux43eux43dux444ux438ux433ux443ux440ux430ux446ux438ux439-ux438-ux441ux43eux441ux442ux43eux44fux43dux438ux439-ux432-redis}

\subsubsection{Обоснование модели
данных}\label{ux43eux431ux43eux441ux43dux43eux432ux430ux43dux438ux435-ux43cux43eux434ux435ux43bux438-ux434ux430ux43dux43dux44bux445}

Проектирование структур данных является фундаментальным этапом,
определяющим производительность, корректность и сопровождаемость системы
в целом. В разрабатываемом программном комплексе Redis выполняет роль
единственного хранилища всех оперативных данных: конфигураций заглушек,
параметров хостов и учётных записей, состояний запущенных процессов,
счётчиков алгоритма Rate Limiter и вспомогательных индексов. Бинарные
файлы артефактов в Redis не хранятся и не кэшируются -- единственной
постоянной копией каждого артефакта является файл на целевом хосте.

Выбор Redis в качестве хранилища обусловлен его архитектурными
характеристиками, подробно рассмотренными в разделе 1.3: время доступа к
любому ключу составляет менее 1 мс, что критически важно для алгоритма
Rate Limiter, выполняемого при обработке каждого входящего запроса.

Redis предоставляет несколько базовых типов данных: строки (String),
хеши (Hash), списки (List), множества (Set) и сортированные множества
(Sorted Set). Для хранения конфигурационных объектов -- заглушек,
хостов, учётных записей -- выбран тип String в сочетании с
JSON-сериализацией. Такой подход позволяет атомарно читать и записывать
всю конфигурацию объекта одной операцией GET/SET, избегая частичных
обновлений и связанных с ними состояний гонки (race conditions). Для
хранения реестров (множеств идентификаторов) используется тип Set,
гарантирующий уникальность элементов. Для счётчиков Rate Limiter
применяется тип String со встроенными атомарными командами INCR и
EXPIRE.

В таблице 2.3 приведена сводная схема всех пространств имён ключей
Redis, используемых в системе.

\emph{Таблица 2.3 -- Пространства имён ключей Redis и используемые типы
данных}

{\def\LTcaptype{none} % do not increment counter
\begin{longtable}[]{@{}
  >{\raggedright\arraybackslash}p{(\linewidth - 6\tabcolsep) * \real{0.2986}}
  >{\raggedright\arraybackslash}p{(\linewidth - 6\tabcolsep) * \real{0.1689}}
  >{\raggedright\arraybackslash}p{(\linewidth - 6\tabcolsep) * \real{0.2893}}
  >{\raggedright\arraybackslash}p{(\linewidth - 6\tabcolsep) * \real{0.2432}}@{}}
\toprule\noalign{}
\begin{minipage}[b]{\linewidth}\centering
\textbf{Пространство имён}
\end{minipage} & \begin{minipage}[b]{\linewidth}\centering
\textbf{Тип Redis}
\end{minipage} & \begin{minipage}[b]{\linewidth}\centering
\textbf{Назначение}
\end{minipage} & \begin{minipage}[b]{\linewidth}\centering
\textbf{Пример ключа}
\end{minipage} \\
\midrule\noalign{}
\endhead
\bottomrule\noalign{}
\endlastfoot
mocks:\{filename\} & String (JSON) & Конфигурация и состояние заглушки &
mocks:payment-stub.jar \\
mocks:registry & Set & Реестр имён всех зарегистрированных заглушек &
mocks:registry \\
hosts:\{hostname\} & String (JSON) & Конфигурация целевого хоста &
hosts:a1b2c3d4-... \\
hosts:registry & Set & Реестр UUID всех зарегистрированных хостов &
hosts:registry \\
accounts:\{uuid\} & String (JSON) & Учётная запись SSH &
accounts:e5f6a7b8-... \\
accounts:registry & Set & Реестр UUID учётных записей &
accounts:registry \\
rate:\{filename\}:\{window\} & String (Integer) & Счётчик Rate Limiter
для временного окна & rate:payment-stub.jar:1700000060 \\
settings:global & String (JSON) & Глобальные параметры приложения &
settings:global \\
\end{longtable}
}

\subsubsection{Структура данных
заглушки}\label{ux441ux442ux440ux443ux43aux442ux443ux440ux430-ux434ux430ux43dux43dux44bux445-ux437ux430ux433ux43bux443ux448ux43aux438}

Центральным объектом системы является запись о заглушке -- имитационном
Java-сервисе, управляемом программным комплексом. Запись хранится по
ключу вида mocks:\{filename\}, где \{filename\} -- имя загруженного
файла артефакта (например, payment-service-stub.jar). Имя файла является
уникальным идентификатором заглушки в системе: при попытке
зарегистрировать файл с именем, уже присутствующим в реестре
mocks:registry, система возвращает ошибку и предлагает оператору
переименовать загружаемый артефакт или удалить существующую запись.

Запись заглушки не хранит путь к бинарному файлу явно. Поскольку файл
всегда располагается в рабочей директории хоста, путь вычисляется
детерминировано при необходимости как \{working\_dir\}/\{filename\}, где
working\_dir берётся из записи хоста, а filename -- из ключа Redis. Это
исключает дублирование данных и единственную точку рассогласования.

Структура JSON-записи заглушки описана в листинге 2.3.

\emph{Листинг 2.3 -- Структура JSON-записи заглушки в Redis}

Ключ: mocks:\{filename\}

Тип: String (JSON)

"payment-service-stub.jar": \{

"hostname": "test-server-01",

"port": 8105,

"pid": 24731,

"status": "RUNNING",

"jvm\_args": "-Xms128m -Xmx256m",

"rate\_limit": 500,

"rate\_limit\_enabled": false,

"registered\_at": "2024-11-01T10:30:00Z",

"started\_at": "2024-11-01T10:35:00Z"

\}

Описание полей записи заглушки приведено в таблице 2.4.

\emph{Таблица 2.4 -- Поля JSON-записи заглушки}

{\def\LTcaptype{none} % do not increment counter
\begin{longtable}[]{@{}
  >{\raggedright\arraybackslash}p{(\linewidth - 6\tabcolsep) * \real{0.2182}}
  >{\raggedright\arraybackslash}p{(\linewidth - 6\tabcolsep) * \real{0.1156}}
  >{\raggedright\arraybackslash}p{(\linewidth - 6\tabcolsep) * \real{0.4337}}
  >{\raggedright\arraybackslash}p{(\linewidth - 6\tabcolsep) * \real{0.2325}}@{}}
\toprule\noalign{}
\begin{minipage}[b]{\linewidth}\centering
\textbf{Поле}
\end{minipage} & \begin{minipage}[b]{\linewidth}\centering
\textbf{Тип}
\end{minipage} & \begin{minipage}[b]{\linewidth}\centering
\textbf{Описание}
\end{minipage} & \begin{minipage}[b]{\linewidth}\centering
\textbf{Значение при регистрации}
\end{minipage} \\
\midrule\noalign{}
\endhead
\bottomrule\noalign{}
\endlastfoot
hostname & String & Hostname целевого хоста; ключ для получения полной
конфигурации хоста из Redis & Выбирается оператором \\
port & Integer & TCP-порт, на котором запущена заглушка на целевом хосте
& null \\
pid & Integer & PID процесса заглушки на целевом хосте & null \\
status & Enum & Текущее состояние: REGISTERED, RUNNING, STOPPED, ERROR &
REGISTERED \\
jvm\_args & String & Аргументы JVM (параметры памяти, системные
свойства) & "" \\
rate\_limit & Integer & Максимально допустимое число запросов в секунду
(RPS) & 0 \\
rate\_limit\_enabled & Boolean & Флаг активности механизма Rate Limiter
& false \\
registered\_at & ISO 8601 & Дата и время регистрации заглушки & Текущее
время UTC \\
started\_at & ISO 8601 & Дата и время последнего запуска процесса &
null \\
\end{longtable}
}

Поля port, pid и started\_at заполняются только после успешного запуска
процесса и сбрасываются в null при его остановке или удалении. Для
получения полного пути к файлу при выполнении SSH-команд Lifecycle
Manager читает поле working\_dir из записи hosts:\{hostname\} и
конкатенирует его с именем файла.

\subsubsection{Структура данных
хоста}\label{ux441ux442ux440ux443ux43aux442ux443ux440ux430-ux434ux430ux43dux43dux44bux445-ux445ux43eux441ux442ux430}

Запись о целевом хосте хранится по ключу вида hosts:\{hostname\}, где
hostname -- сетевое имя или IP-адрес хоста. Hostname одновременно
выполняет роль уникального идентификатора записи и адреса для
SSH/SCP-подключения, что исключает необходимость в отдельном поле
адреса. Прямой доступ к конфигурации хоста по известному hostname
осуществляется за одну операцию чтения без дополнительного индекса.

Структура JSON-записи хоста описана в листинге 2.4.

\emph{Листинг 2.4 -- Структура JSON-записи хоста в Redis}

Ключ: hosts:\{hostname\}

Тип: String (JSON)

"hosts:test-server-01": \{

"ssh\_port": 22,

"account\_uuid": "e5f6a7b8-c9d0-1234-efab-567890abcdef",

"working\_dir": "/opt/mock-services",

"java\_path": "/usr/lib/jvm/java-17-openjdk-amd64/bin/java",

"mock\_port\_min": 8100,

"mock\_port\_max": 8200,

"description": "Стенд нагрузочного тестирования -- контур препром.",

"status": "AVAILABLE",

"last\_checked\_at":"2024-11-01T10:29:55Z"

\}

Поле status отражает результат последней проверки доступности
SSH-соединения, выполняемой фоновой задачей с настраиваемым интервалом.
Допустимые значения: AVAILABLE (соединение успешно установлено),
UNAVAILABLE (попытка соединения завершилась ошибкой) и UNKNOWN (проверка
ещё не выполнялась). Поле last\_checked\_at фиксирует время последней
проверки для отображения актуальности статуса в веб-интерфейсе.

Поле account\_uuid содержит UUID учётной записи SSH из реестра accounts,
используемой для подключения к данному хосту. Lifecycle Manager читает
это поле при каждой SSH/SCP-операции для получения соответствующей
записи accounts:\{uuid\} с зашифрованным паролем.

\subsubsection{Структура данных учётной
записи}\label{ux441ux442ux440ux443ux43aux442ux443ux440ux430-ux434ux430ux43dux43dux44bux445-ux443ux447ux451ux442ux43dux43eux439-ux437ux430ux43fux438ux441ux438}

Запись об учётной записи ОС Linux хранится по ключу вида
accounts:\{uuid\}. Учётная запись содержит данные, необходимые для
аутентификации при SSH/SCP-операциях на целевых хостах. Пароль хранится
в зашифрованном виде -- компонент Crypto Service шифрует его алгоритмом
AES-128 (Fernet) перед записью в Redis. Ключ шифрования хранится на
файловой системе сервера в файле с правами 600.

Структура JSON-записи учётной записи описана в листинге 2.5.

\emph{Листинг 2.5 -- Структура JSON-записи учётной записи в Redis}

Ключ: accounts:\{uuid\}

Тип: String (JSON)

"accounts:e5f6a7b8-c9d0-1234-efab-567890abcdef": \{

"username": "mock-runner",

"password\_enc": "gAAAAABl...base64-encoded-ciphertext...",

"description": "Сервисный пользователь для запуска заглушек",

"created\_at": "2024-10-15T09:00:00Z"

\}

Поле password\_enc содержит Base64-кодированный зашифрованный блок,
полученный библиотекой cryptography (Fernet). При необходимости
установить SSH-соединение Crypto Service читает это поле и выполняет
дешифрование в оперативной памяти процесса; расшифрованный пароль
никогда не записывается на диск и не сохраняется в Redis. Такой подход
обеспечивает защиту паролей в случае несанкционированного получения
дампа данных Redis.

\subsubsection{Структуры данных
реестров}\label{ux441ux442ux440ux443ux43aux442ux443ux440ux44b-ux434ux430ux43dux43dux44bux445-ux440ux435ux435ux441ux442ux440ux43eux432}

Для хранения множеств идентификаторов всех зарегистрированных объектов
используется тип Set. Каждый реестр хранится под фиксированным ключом и
содержит строковые идентификаторы соответствующих объектов.

\textbf{Реестр заглушек (mocks:registry)} -- множество имён файлов всех
зарегистрированных заглушек. При регистрации новой заглушки её имя
добавляется командой SADD, при удалении -- удаляется командой SREM.
Команда SMEMBERS возвращает полный список для формирования Dashboard.

\textbf{Реестр хостов (hosts:registry)} -- множество
hostname-идентификаторов зарегистрированных целевых хостов (сетевое имя
или IP-адрес). При регистрации нового хоста его hostname добавляется
командой SADD, при удалении -- командой SREM. Команда SMEMBERS
возвращает полный список для отображения в интерфейсе настроек.

\textbf{Реестр учётных записей (accounts:registry)} -- множество UUID
зарегистрированных учётных записей SSH.

Применение типа Set гарантирует уникальность идентификаторов: повторное
добавление уже существующего элемента командой SADD не создаёт дубликата
и возвращает 0, что используется системой для обнаружения попытки
регистрации заглушки с дублирующим именем файла.

Схема взаимосвязей между реестрами и конфигурационными записями
приведена на рисунке 2.6.

\includesvg[width=6.01042in,height=5.26042in]{media/image10.svg}Рисунок
2.6 -- Схема взаимосвязей объектов в Redis

@startuml redis\_relations

skinparam defaultFontName Arial

skinparam backgroundColor \#FFFFFF

package "Redis" \{

rectangle "mocks:registry\textbackslash n(Set)" as MR \#D5E8D4

rectangle "mocks:\{filename\}\textbackslash n(String/JSON)" as M
\#DAE8FC

rectangle "hosts:registry\textbackslash n(Set)" as HR \#D5E8D4

rectangle "hosts:\{uuid\}\textbackslash n(String/JSON)" as H \#DAE8FC

rectangle "accounts:registry\textbackslash n(Set)" as AR \#D5E8D4

rectangle "accounts:\{uuid\}\textbackslash n(String/JSON)" as A \#DAE8FC

rectangle "rate:\{filename\}:\{window\}\textbackslash n(String/Integer)"
as RL \#FFF2CC

rectangle "settings:global\textbackslash n(String/JSON)" as SG \#F8CECC

MR -\/-\textgreater{} M : SMEMBERS → GET

M -\/-\textgreater{} H : host\_id → GET

H -\/-\textgreater{} A : account\_id → GET

HR -\/-\textgreater{} H : SMEMBERS → GET

AR -\/-\textgreater{} A : SMEMBERS → GET

M -\/-\textgreater{} RL : filename → INCR/GET

\}

@enduml

\subsubsection{Структура данных счётчика Rate
Limiter}\label{ux441ux442ux440ux443ux43aux442ux443ux440ux430-ux434ux430ux43dux43dux44bux445-ux441ux447ux451ux442ux447ux438ux43aux430-rate-limiter}

Алгоритм Rate Limiter реализован на основе паттерна «Счётчик с
фиксированным временным окном» (Fixed Window Counter). Для каждой
заглушки и каждого временного окна в Redis хранится отдельный
ключ-счётчик. Структура ключа: rate:\{filename\}:\{window\}, где
\{window\} -- метка временного окна, вычисляемая как целочисленное
деление текущего Unix-timestamp на длину окна в секундах.

Например, при длине окна 1 секунда и текущем времени 1700000060.7 метка
окна равна 1700000060, а ключ для заглушки payment-service-stub.jar
примет вид rate:payment-service-stub.jar:1700000060. Все запросы,
поступившие в течение одной секунды, инкрементируют один и тот же
счётчик.

Алгоритм обработки каждого запроса через Rate Limiter описан в листинге
2.7.

\emph{Листинг 2.7 -- Псевдокод алгоритма Fixed Window Counter}

function check\_rate\_limit(filename: str, limit: int, window\_size: int
= 1) -\textgreater{} bool:

\# Шаг 1: вычислить метку текущего окна

window = floor(current\_unix\_timestamp() / window\_size)

key = f"rate:\{filename\}:\{window\}"

\# Шаг 2: атомарный инкремент счётчика

count = INCR(key)

\# Шаг 3: установить TTL при создании нового окна

if count == 1:

EXPIRE(key, window\_size * 2) \# TTL = 2 окна (запас на сдвиг времени)

\# Шаг 4: проверить лимит

if count \textgreater{} limit:

return False \# запрос отклонён → HTTP 429

return True \# запрос разрешён

Установка TTL через команду EXPIRE выполняется только при создании
нового счётчика (когда INCR возвращает 1), то есть при первом запросе в
новом временном окне. TTL устанавливается равным удвоенному размеру окна
для гарантированного самоочищения устаревших счётчиков даже при
незначительных расхождениях системного времени между клиентом и
сервером. Это исключает накопление «мёртвых» ключей в Redis и
предотвращает утечку памяти при длительной эксплуатации системы.

Атомарность операции INCR является ключевым свойством данного подхода:
Redis выполняет её как единую неделимую операцию, что гарантирует
корректный подсчёт при конкурентном поступлении тысяч запросов в рамках
одного временного окна без использования транзакций или блокировок.

\subsubsection{Глобальные настройки
системы}\label{ux433ux43bux43eux431ux430ux43bux44cux43dux44bux435-ux43dux430ux441ux442ux440ux43eux439ux43aux438-ux441ux438ux441ux442ux435ux43cux44b}

Глобальные конфигурационные параметры приложения, не привязанные к
конкретному хосту или заглушке, хранятся под ключом settings:global.
Запись содержит параметры, задаваемые при первоначальной настройке
системы и редко изменяемые в процессе эксплуатации.

Структура JSON-записи глобальных настроек описана в листинге 2.8.

\emph{Листинг 2.8 -- Структура глобальных настроек в Redis}

Ключ: settings:global

Тип: String (JSON)

"settings:global": \{

"rate\_limit\_window\_size": 1,

"host\_check\_interval\_sec": 30,

"proxy\_timeout\_sec": 10,

"log\_retention\_lines": 1000

\}

Описание полей глобальных настроек приведено в таблице 2.5.

\emph{Таблица 2.5 -- Поля глобальных настроек системы}

{\def\LTcaptype{none} % do not increment counter
\begin{longtable}[]{@{}
  >{\raggedright\arraybackslash}p{(\linewidth - 6\tabcolsep) * \real{0.2781}}
  >{\raggedright\arraybackslash}p{(\linewidth - 6\tabcolsep) * \real{0.0984}}
  >{\raggedright\arraybackslash}p{(\linewidth - 6\tabcolsep) * \real{0.1653}}
  >{\raggedright\arraybackslash}p{(\linewidth - 6\tabcolsep) * \real{0.4582}}@{}}
\toprule\noalign{}
\begin{minipage}[b]{\linewidth}\centering
Параметр
\end{minipage} & \begin{minipage}[b]{\linewidth}\centering
Тип
\end{minipage} & \begin{minipage}[b]{\linewidth}\centering
По умолчанию
\end{minipage} & \begin{minipage}[b]{\linewidth}\centering
Описание
\end{minipage} \\
\midrule\noalign{}
\endhead
\bottomrule\noalign{}
\endlastfoot
rate\_limit\_window\_size & Integer & 1 & Длина временного окна
алгоритма Rate Limiter в секундах \\
host\_check\_interval\_sec & Integer & 30 & Интервал фоновой проверки
доступности хостов (SSH-пинг) в секундах \\
proxy\_timeout\_sec & Integer & 10 & Тайм-аут ожидания ответа от целевой
заглушки при проксировании (секунды) \\
log\_retention\_lines & Integer & 1000 & Максимальное число строк лога
процесса, хранимых в памяти системы \\
\end{longtable}
}

\subsubsection{Персистентность данных и стратегия
восстановления}\label{ux43fux435ux440ux441ux438ux441ux442ux435ux43dux442ux43dux43eux441ux442ux44c-ux434ux430ux43dux43dux44bux445-ux438-ux441ux442ux440ux430ux442ux435ux433ux438ux44f-ux432ux43eux441ux441ux442ux430ux43dux43eux432ux43bux435ux43dux438ux44f}

Redis настроен в режиме AOF (Append Only File) с параметром appendfsync
everysec: каждая операция записи фиксируется в журнале, синхронизируемом
с диском раз в секунду. Максимальная потеря данных при аварии составляет
не более одной секунды.

Стратегия восстановления после перезапуска: управляющее приложение
читает реестры из Redis и для каждой заглушки со статусом RUNNING
проверяет фактическое наличие процесса на целевом хосте через SSH. Если
процесс обнаружен -- статус подтверждается. Если нет -- статус
обновляется до STOPPED, поля port, pid и started\_at сбрасываются.
Бинарный файл на целевом хосте не затрагивается и остаётся в рабочей
директории, готовый к повторному запуску.

В таблице 2.6 представлена сводная матрица операций чтения и записи
Redis по компонентам системы.

\emph{Таблица 2.6 -- Матрица операций Redis по компонентам системы}

{\def\LTcaptype{none} % do not increment counter
\begin{longtable}[]{@{}
  >{\raggedright\arraybackslash}p{(\linewidth - 6\tabcolsep) * \real{0.1653}}
  >{\raggedright\arraybackslash}p{(\linewidth - 6\tabcolsep) * \real{0.3189}}
  >{\raggedright\arraybackslash}p{(\linewidth - 6\tabcolsep) * \real{0.3162}}
  >{\raggedright\arraybackslash}p{(\linewidth - 6\tabcolsep) * \real{0.1996}}@{}}
\toprule\noalign{}
\begin{minipage}[b]{\linewidth}\centering
\textbf{Компонент}
\end{minipage} & \begin{minipage}[b]{\linewidth}\centering
\textbf{Операция}
\end{minipage} & \begin{minipage}[b]{\linewidth}\centering
\textbf{Ключ / пространство имён}
\end{minipage} & \begin{minipage}[b]{\linewidth}\centering
\textbf{Redis-команды}
\end{minipage} \\
\midrule\noalign{}
\endhead
\bottomrule\noalign{}
\endlastfoot
Proxy Engine & Чтение конфигурации заглушки & mocks:\{filename\} &
GET \\
Rate Limiter & Инкремент и проверка счётчика &
rate:\{filename\}:\{window\} & INCR, EXPIRE \\
Lifecycle Manager & Регистрация заглушки (после SCP) &
mocks:\{filename\}, mocks:registry & SET, SADD \\
Lifecycle Manager & Обновление статуса/порта/PID & mocks:\{filename\} &
GET → modify → SET \\
Lifecycle Manager & Удаление заглушки (после SSH rm) &
mocks:\{filename\}, mocks:registry & DEL, SREM \\
Lifecycle Manager & Получение параметров хоста & hosts:\{hostname\} &
GET \\
Lifecycle Manager & Получение учётной записи &
accounts:\{account\_name\} & GET \\
Settings Module & Сохранение/обновление хоста & hosts:\{hostname\},
hosts:registry & SET, SADD \\
Settings Module & Сохранение/обновление УЗ & accounts:\{name\},
accounts:registry & SET, SADD \\
Settings Module & Переключение rate\_limit\_enabled & mocks:\{filename\}
& GET → modify → SET \\
Metrics Exporter & Чтение всех заглушек для метрик & mocks:registry +
mocks:\{filename\} & SMEMBERS, GET (pipeline) \\
Background Task & Обновление статуса доступности хоста &
hosts:\{hostname\} & GET → modify → SET \\
\end{longtable}
}

Для операции Metrics Exporter применяется механизм конвейеризации команд
Redis (Pipeline): сначала читается полный список ключей из реестра,
затем все операции чтения конфигураций отправляются единым пакетом. Это
снижает суммарные сетевые задержки при агрегации данных о большом числе
заглушек с O(n) до фактически O(1) по числу сетевых round-trip.

\end{document}
